\documentclass[10pt,aspectratio=169,t]{beamer}

% ============================================================================
% GÓI NGÔN NGỮ TIẾNG VIỆT & BẢNG MÃ FONT
% ============================================================================
\usepackage[utf8]{inputenc}
\usepackage[T5]{fontenc}
\usepackage[vietnamese]{babel}
\shorthandoff{"}
\AtBeginDocument{\shorthandoff{"}}

% ============================================================================
% CÁC GÓI TOÁN HỌC & ĐỒ HỌA TIKZ
% ============================================================================
\usepackage{amsmath, amssymb, amsfonts, mathtools, bm}
\usepackage{booktabs, tabularx, array, multirow}
\usepackage{tikz}
\usetikzlibrary{arrows.meta, positioning, shapes.geometric, calc, fit, backgrounds}

% Toán tử tuỳ chỉnh
\DeclareMathOperator*{\argmax}{arg\,max}
\DeclareMathOperator*{\argmin}{arg\,min}
\DeclareMathOperator{\E}{\mathbb{E}}
\DeclareMathOperator{\DKL}{D_{\mathrm{KL}}}
\DeclareMathOperator{\clip}{\mathrm{clip}}
\DeclareMathOperator{\mean}{\mathrm{mean}}
\DeclareMathOperator{\std}{\mathrm{std}}
\newcommand{\cancel}[1]{\textcolor{stanfordred}{\underbrace{#1}_{\text{Triệt tiêu } = 0}}}

% ============================================================================
% CẤU HÌNH GIAO DIỆN BEAMER HỌC THUẬT (16:9)
% ============================================================================
\usetheme{Madrid}
\usecolortheme{default}
\setbeamertemplate{navigation symbols}{}

\definecolor{stanfordred}{RGB}{140, 21, 21}
\definecolor{navyblue}{RGB}{15, 45, 90}
\definecolor{tealgreen}{RGB}{0, 110, 90}
\definecolor{amberorange}{RGB}{195, 95, 0}
\definecolor{softblue}{RGB}{238, 245, 253}
\definecolor{softred}{RGB}{253, 240, 240}
\definecolor{softgreen}{RGB}{236, 248, 241}
\definecolor{softyellow}{RGB}{254, 250, 235}
\definecolor{slate}{RGB}{70, 80, 95}

\setbeamercolor{structure}{fg=navyblue}
\setbeamercolor{palette primary}{bg=navyblue, fg=white}
\setbeamercolor{palette secondary}{bg=stanfordred, fg=white}
\setbeamercolor{palette tertiary}{bg=navyblue!90!black, fg=white}
\setbeamercolor{frametitle}{bg=navyblue, fg=white}
\setbeamercolor{framesubtitle}{fg=softyellow!95!white}
\setbeamercolor{title}{bg=navyblue, fg=white}

\setbeamercolor{block title}{bg=navyblue, fg=white}
\setbeamercolor{block body}{bg=softblue, fg=black}

\setbeamercolor{block title alerted}{bg=stanfordred, fg=white}
\setbeamercolor{block body alerted}{bg=softred, fg=black}

\setbeamercolor{block title example}{bg=tealgreen, fg=white}
\setbeamercolor{block body example}{bg=softgreen, fg=black}

\setbeamertemplate{itemize items}[circle]
\setbeamertemplate{enumerate items}[default]
\setbeamertemplate{section in toc}[sections numbered]
\setbeamertemplate{subsection in toc}[subsections numbered]

\setbeamerfont{frametitle}{size=\normalsize,series=\bfseries}
\setbeamerfont{framesubtitle}{size=\footnotesize,series=\mdseries}

% Tự động chèn Slide Mục lục Chuyển Phần ở đầu mỗi Phần lớn
\AtBeginSection[]{
  \begin{frame}{Mục lục Tiến trình Seminar (2020--2026)}{Cấu trúc Học thuật \& Vị trí Phần Hiện tại}
    \scriptsize
    \tableofcontents[currentsection, subsectionstyle=show/show/hide]
  \end{frame}
}

% Định nghĩa tiêu đề phân mục phẳng
\newcommand{\subhead}[1]{%
  \vspace{0.10cm}%
  {\normalsize\bfseries\textcolor{navyblue}{#1}}\par\vspace{0.04cm}%
  {\color{navyblue!40}\hrule height 0.6pt}\vspace{0.10cm}%
}
\newcommand{\limithead}[1]{%
  \vspace{0.10cm}%
  {\normalsize\bfseries\textcolor{stanfordred}{#1}}\par\vspace{0.04cm}%
  {\color{stanfordred!45}\hrule height 0.6pt}\vspace{0.10cm}%
}
\newcommand{\motivehead}[1]{%
  \vspace{0.10cm}%
  {\normalsize\bfseries\textcolor{tealgreen}{#1}}\par\vspace{0.04cm}%
  {\color{tealgreen!45}\hrule height 0.6pt}\vspace{0.10cm}%
}

\AtBeginDocument{%
  \setlength{\abovedisplayskip}{5pt}%
  \setlength{\belowdisplayskip}{5pt}%
  \setlength{\abovedisplayshortskip}{3pt}%
  \setlength{\belowdisplayshortskip}{3pt}%
}

% Footer học thuật
\setbeamertemplate{footline}{
  \leavevmode%
  \hbox{%
  \begin{beamercolorbox}[wd=.26\paperwidth,ht=2.4ex,dp=1ex,left]{author in head/foot}%
    \hspace*{2ex}\usebeamerfont{author in head/foot}\textbf{Seminar Học thuật: Post-Training (2020--2026)}
  \end{beamercolorbox}%
  \begin{beamercolorbox}[wd=.56\paperwidth,ht=2.4ex,dp=1ex,center]{title in head/foot}%
    \usebeamerfont{title in head/foot}\insertsectionhead
  \end{beamercolorbox}%
  \begin{beamercolorbox}[wd=.18\paperwidth,ht=2.4ex,dp=1ex,right]{date in head/foot}%
    \textbf{Trang \insertframenumber{} / \inserttotalframenumber}\hspace*{2ex}
  \end{beamercolorbox}}%
  \vskip0pt%
}

% ============================================================================
% ĐỊNH NGHĨA STYLE TIKZ CHO CÁC SƠ ĐỒ LUỒNG DỮ LIỆU
% ============================================================================
\tikzset{
  >=Stealth,
  data/.style={
    rectangle, rounded corners=4pt, draw=amberorange!85!black, line width=0.9pt,
    fill=softyellow, text width=3.0cm, align=center, minimum height=0.95cm, inner sep=4pt, font=\scriptsize\sffamily
  },
  trainmodel/.style={
    rectangle, rounded corners=4pt, draw=navyblue, line width=1.1pt,
    fill=softblue, text width=3.1cm, align=center, minimum height=1.05cm, inner sep=4pt, font=\scriptsize\sffamily\bfseries
  },
  frozenmodel/.style={
    rectangle, rounded corners=4pt, draw=slate, line width=0.9pt, dashed,
    fill=gray!12, text width=3.1cm, align=center, minimum height=1.05cm, inner sep=4pt, font=\scriptsize\sffamily
  },
  rewardbox/.style={
    rectangle, rounded corners=4pt, draw=tealgreen, line width=1.0pt,
    fill=softgreen, text width=3.0cm, align=center, minimum height=0.95cm, inner sep=4pt, font=\scriptsize\sffamily
  },
  lossbox/.style={
    rectangle, rounded corners=4pt, draw=stanfordred, line width=1.1pt,
    fill=softred, text width=3.2cm, align=center, minimum height=1.0cm, inner sep=4pt, font=\scriptsize\sffamily\bfseries
  },
  flow/.style={->, line width=1.05pt, draw=navyblue!85!black, rounded corners=3pt},
  gradflow/.style={->, line width=1.1pt, draw=stanfordred, dashed, rounded corners=3pt}
}

% ============================================================================
% THÔNG TIN BÀI SEMINAR KHOA HỌC
% ============================================================================
\title[Seminar: Post-Training cho LLMs]{CƠ SỞ TOÁN HỌC VÀ SỰ TIẾN HÓA THUẬT TOÁN HẬU HUẤN LUYỆN MÔ HÌNH NGÔN NGỮ LỚN (2020--2026)}
\subtitle{Phần 1: SFT (2020--2021) $\cdot$ Phần 2: RLHF \& PPO (2022) $\cdot$ Phần 3: DPO (2023--2024) $\cdot$ Phần 4: RLVR \& GRPO (2024--2026) $\cdot$ Phần 5: Chưng cất Tri thức}
\author[Báo cáo Seminar Học thuật]{Chuyên đề Seminar Khoa học Máy tính \& Trí tuệ Nhân tạo}
\institute[]{Giải tích Tối ưu hóa, Kiến trúc Luồng Dữ liệu và Động cơ Tiến hóa Phương pháp}
\date{Thời lượng trình bày: 90--120 Phút}

\begin{document}

\begin{frame}[plain]
  \vspace{0.2cm}
  \titlepage
\end{frame}

\begin{frame}{Mục lục Tổng quan Báo cáo Seminar}{Tiến trình Lịch sử (2020--2026) \& Chưng cất Tri thức}
  \scriptsize
  \tableofcontents[hideallsubsections]
  \vspace{0.15cm}
  \begin{exampleblock}{Cấu trúc Tổng kết Cuối mỗi Phần}
    \scriptsize
    Kết thúc mỗi Phần đều tổng hợp rõ: \textbf{(1) Đóng góp toán học \& tối ưu hệ thống đạt được so với giai đoạn trước}, và \textbf{(2) Giới hạn tồn đọng chuyển giao cho Phần tiếp theo giải quyết}.
  \end{exampleblock}
\end{frame}

% ============================================================================
% PHẦN 1 (2020 - 2021): TỪ TIỀN HUẤN LUYỆN ĐẾN TINH CHỈNH CÓ GIÁM SÁT (SFT)
% ============================================================================
\section{Phần 1 (2020--2021): Từ Tiền huấn luyện đến Tinh chỉnh có Giám sát (SFT)}

\subsection{1.1. Bối cảnh Lịch sử (2020--2026) \& Sự Lệch Mục tiêu của Tiền huấn luyện}

% ----------------------------------------------------------------------------
% SLIDE: MÔ HÌNH TIỀN HUẤN LUYỆN (GPT-3, 2020)
% ----------------------------------------------------------------------------
\begin{frame}{Phần 1.1: Mô hình Ngôn ngữ Tiền huấn luyện (GPT-3 -- 2020)}{Hàm Mục tiêu Tự hồi quy và Cơ chế Nén Tri thức Nền tảng (Brown et al., 2020)}
  \begin{block}{Công thức 1: Hàm Mục tiêu Dự đoán Token Tiếp theo (Autoregressive MLE)}
    \[
      \mathcal{L}_{\text{Pretrain}}(\theta) = -\sum_{t=1}^{T} \log P_\theta(u_t \mid u_1, u_2, \dots, u_{t-1})
    \]
  \end{block}

  \subhead{Đặc tính Toán học của Mô hình Nền (Base Model)}
  \small
  \begin{itemize}\setlength{\itemsep}{7pt}
    \item \textbf{Xấp xỉ phân phối biên ngôn ngữ:} Học phân phối có điều kiện $P_\theta(u_t \mid u_{<t})$ trên kho văn bản quy mô $10^{12}\text{--}10^{13}$ token.
    \item \textbf{Nội hóa tri thức thế giới:} Nén cấu trúc cú pháp, tri thức sự kiện (factual knowledge) và quy luật suy diễn vào không gian tham số $\theta$.
    \item \textbf{Giới hạn năm 2020:} Vận hành thuần túy như bộ nối tiếp chuỗi xác suất; phụ thuộc hoàn toàn vào kỹ thuật đặt ví dụ trong ngữ cảnh (\textit{Few-shot In-Context Learning}).
  \end{itemize}
\end{frame}

% ----------------------------------------------------------------------------
% SLIDE: SỰ LỆCH MỤC TIÊU (OBJECTIVE MISMATCH)
% ----------------------------------------------------------------------------
\begin{frame}{Phần 1.1: Sự Lệch Mục tiêu (Objective Mismatch)}{Bất tương thích giữa Huấn luyện Dự đoán Văn bản và Tác vụ Giải quyết Yêu cầu}
  \begin{columns}[T,onlytextwidth]
    \begin{column}{0.48\textwidth}
      \limithead{Hạn chế của Tiền huấn luyện (Pre-training)}
      \small
      \begin{itemize}\setlength{\itemsep}{7pt}
        \item \textbf{Lệch mục tiêu tối ưu hóa:}
          \begin{itemize}\footnotesize\setlength{\itemsep}{3pt}
            \item Đầu vào $x$: \textit{``Giải thích nguyên lý của laser?''}
            \item Đầu ra Base Model: Sinh tiếp các câu hỏi số 2, số 3 thay vì trả lời.
          \end{itemize}
        \item \textbf{Không phân biệt chất lượng dữ liệu:}
          \begin{itemize}\footnotesize\setlength{\itemsep}{3pt}
            \item Tái tạo cả sai số thực tế, định kiến và nội dung độc hại từ tập dữ liệu web.
          \end{itemize}
      \end{itemize}
    \end{column}
    \begin{column}{0.48\textwidth}
      \motivehead{Mục tiêu của Hậu huấn luyện (Post-Training)}
      \small
      \begin{itemize}\setlength{\itemsep}{7pt}
        \item \textbf{Định hướng phân phối phản hồi:} Chuyển dịch từ $P_\theta(u_t \mid u_{<t})$ sang phân phối phản hồi theo chỉ thị $\pi_\theta(y \mid x)$.
        \item \textbf{Ba tiêu chuẩn căn chỉnh (3H Criteria):}
          \begin{enumerate}\footnotesize\setlength{\itemsep}{3pt}
            \item \textbf{Helpful (Hữu ích):} Giải quyết đúng yêu cầu.
            \item \textbf{Honest (Trung thực):} Hạn chế bịa đặt thông tin.
            \item \textbf{Harmless (An toàn):} Từ chối tác vụ độc hại.
          \end{enumerate}
      \end{itemize}
    \end{column}
  \end{columns}
\end{frame}

% ----------------------------------------------------------------------------
% SLIDE: BIÊN NIÊN SỬ TIẾN HÓA POST-TRAINING 2020 - 2026 (TOÀN TRANG)
% ----------------------------------------------------------------------------
\begin{frame}{Phần 1.1: Biên niên sử Tiến hóa Hậu Huấn luyện (2020 -- 2026)}{Dòng Chảy Lịch sử qua Năm Giai đoạn Chuyển dịch Thuật toán}
  \vspace{0.12cm}
  \begin{center}
    \resizebox{0.97\textwidth}{!}{%
      \begin{tikzpicture}[
        era/.style={rectangle, rounded corners=4pt, draw=navyblue, line width=1.0pt, fill=softblue, text width=3.1cm, minimum height=1.25cm, align=center, font=\scriptsize\sffamily\bfseries},
        detail/.style={rectangle, rounded corners=4pt, draw=tealgreen!85!black, line width=0.85pt, fill=softgreen, text width=3.1cm, minimum height=2.15cm, align=left, inner sep=4pt, font=\tiny\sffamily}
      ]
        \node[era] (e1) at (0, 1.35) {Giai đoạn 1 (2020--2021)\\\textcolor{stanfordred}{Phần 1: SFT}\\(FLAN, T0)};
        \node[era] (e2) at (3.65, 1.35) {Giai đoạn 2 (2022)\\\textcolor{stanfordred}{Phần 2: RLHF + PPO}\\(InstructGPT, ChatGPT)};
        \node[era] (e3) at (7.3, 1.35) {Giai đoạn 3 (2023--2024)\\\textcolor{stanfordred}{Phần 3: DPO \& Biến thể}\\(Llama 2/3, SimPO)};
        \node[era, fill=softyellow, draw=amberorange!90!black] (e4) at (10.95, 1.35) {Giai đoạn 4 (2024--2026)\\\textcolor{stanfordred}{Phần 4: RLVR + GRPO}\\(o1/o3, DeepSeek-R1)};
        \node[era, fill=softred, draw=stanfordred] (e5) at (14.6, 1.35) {Chuyên đề Cuối\\\textcolor{stanfordred}{Phần 5: Chưng cất}\\(Distillation \& Merging)};

        \node[detail] (d1) at (0, -0.75) {
          \textbf{Đóng góp:} Học có giám sát trên mẫu chuẩn $\mathcal{L}_{\text{SFT}}(y^* \mid x)$.\\[4pt]
          \textbf{Hạn chế:} \textit{Exposure Bias}, ảo giác, thiếu tín hiệu từ mẫu âm.
        };
        \node[detail] (d2) at (3.65, -0.75) {
          \textbf{Đóng góp:} Hàm thưởng $r_\phi(y_w \succ y_l)$ + Điều chuẩn $\beta\DKL$ + PPO 4 mạng.\\[4pt]
          \textbf{Hạn chế:} Quá nặng VRAM (4 mô hình), đạo hàm RL bất ổn.
        };
        \node[detail] (d3) at (7.3, -0.75) {
          \textbf{Đóng góp:} Nghiệm đóng khử $Z(x)$, bỏ $r_\phi$ \& $V_\psi$, chỉ cần 2 mô hình.\\[4pt]
          \textbf{Hạn chế:} Học ngoại tuyến (\textit{Off-policy}), bị \textit{Reward Hacking} ở Toán/Code.
        };
        \node[detail, fill=softyellow!70!white, draw=amberorange!90!black] (d4) at (10.95, -0.75) {
          \textbf{Đóng góp:} Kiểm chứng tất định + Chuẩn hóa nhóm $G$ bỏ Critic $V_\psi$ $\to$ Chuỗi \texttt{<think>}.\\[4pt]
          \textbf{Hạn chế:} Mô hình nhỏ bị $\nabla_\theta = \mathbf{0}$.
        };
        \node[detail, fill=softred!60!white, draw=stanfordred] (d5) at (14.6, -0.75) {
          \textbf{Đóng góp:} Chuyển giao phân phối $P_T^{(\tau)}$ và chuỗi \texttt{<think>} từ Thầy sang Trò.\\[4pt]
          \textbf{Mở rộng:} Reverse KL, LoRA, Merging \& Đánh giá.
        };

        \draw[flow] (e1.east) -- (e2.west);
        \draw[flow] (e2.east) -- (e3.west);
        \draw[flow] (e3.east) -- (e4.west);
        \draw[flow] (e4.east) -- (e5.west);

        \draw[flow] (e1.south) -- (d1.north);
        \draw[flow] (e2.south) -- (d2.north);
        \draw[flow] (e3.south) -- (d3.north);
        \draw[flow] (e4.south) -- (d4.north);
        \draw[flow] (e5.south) -- (d5.north);
      \end{tikzpicture}%
    }
  \end{center}
  \vspace{0.04cm}
  \footnotesize
  \textbf{Quy luật tiến hóa:} Giới hạn toán học và hệ thống của giai đoạn trước là động cơ trực tiếp sinh ra thuật toán của giai đoạn kế tiếp.
\end{frame}

\subsection{1.2. Tinh chỉnh theo Chỉ thị có Giám sát (SFT -- Giai đoạn 2021--2022)}

% ----------------------------------------------------------------------------
% SLIDE: SUPERVISED FINE-TUNING (SFT) & KHUÔN MẪU HỘI THOẠI
% ----------------------------------------------------------------------------
\begin{frame}{Phần 1.2: Tinh chỉnh có Giám sát (SFT) --- Cấu trúc Hóa Chuỗi Tương tác}{Phương pháp Học Bắt chước Chuyên gia (Behavioral Cloning -- Wei et al., 2022)}
  \subhead{Định dạng Dữ liệu Giám sát và Token Điều khiển}
  \small
  \begin{itemize}\setlength{\itemsep}{6pt}
    \item \textbf{Tập dữ liệu:} $\mathcal{D}_{\text{SFT}} = \{(x^{(i)}, y^{*(i)})\}_{i=1}^{N}$ gồm truy vấn đầu vào $x$ và mẫu phản hồi chuẩn $y^*$.
    \item \textbf{Phân vùng ngữ cảnh bằng Token đặc biệt (Control Tokens):}
      \begin{itemize}\footnotesize\setlength{\itemsep}{3pt}
        \item \texttt{<|im\_start|>system ... <|im\_end|>} $\longrightarrow$ Ràng buộc hệ thống toàn cục.
        \item \texttt{<|im\_start|>user ... <|im\_end|>} $\longrightarrow$ Miền truy vấn đầu vào $x$.
        \item \texttt{<|im\_start|>assistant ... <|im\_end|>} $\longrightarrow$ Miền sinh phản hồi mục tiêu $y^*$.
      \end{itemize}
  \end{itemize}

  \vspace{0.1cm}
  \begin{exampleblock}{Ý nghĩa Toán học của Token Kết thúc \texttt{<EOS>} (\texttt{<|im\_end|>})}
    \footnotesize
    Đưa \texttt{<EOS>} vào miền tính đạo hàm giúp mô hình ước lượng chính xác xác suất dừng chuỗi $P_\theta(\texttt{<EOS>} \mid x, y_{\le T}^*)$, triệt tiêu hiện tượng sinh văn bản tràn lan của Base Model.
  \end{exampleblock}
\end{frame}

% ----------------------------------------------------------------------------
% SLIDE: SƠ ĐỒ LUỒNG DỮ LIỆU HUẤN LUYỆN SFT & PROMPT MASKING
% ----------------------------------------------------------------------------
\begin{frame}{Phần 1.2: Sơ đồ Luồng Dữ liệu Huấn luyện SFT \& Mặt nạ Hàm Mất mát}{Cơ chế Truyền Thuận Toàn chuỗi và Lan truyền Ngược có Chọn lọc (Prompt Masking)}
  \vspace{0.15cm}
  \begin{center}
    \resizebox{0.96\textwidth}{!}{%
      \begin{tikzpicture}
        \node[data, text width=3.6cm, fill=gray!18, draw=slate] (prompt) at (0, 1.0) {
          \textbf{Truy vấn Đầu vào $x$}\\
          \texttt{<|user|> Nội dung câu hỏi...}\\
          Mặt nạ: $\mathbf{m_t = 0}$ (Triệt tiêu lỗi)
        };
        \node[data, text width=3.6cm, fill=softgreen, draw=tealgreen] (resp) at (0, -1.0) {
          \textbf{Mẫu Phản hồi Chuẩn $y^*$}\\
          \texttt{<|assistant|> Lời giải... <|end|>}\\
          Mặt nạ: $\mathbf{m_t = 1}$ (Tính lỗi)
        };

        \node[trainmodel, text width=3.4cm, minimum height=2.6cm] (llm) at (4.7, 0) {
          \textbf{Mô hình Chính sách $\pi_\theta$}\\
          (Đang tối ưu hóa)\\[6pt]
          \normalfont\scriptsize Biểu diễn ngữ cảnh nối ghép $[x \oplus y^*]$ qua Causal Self-Attention
        };

        \node[rewardbox, text width=3.3cm, minimum height=1.5cm] (prob) at (9.2, 0) {
          \textbf{Phân phối Xác suất}\\
          $\pi_\theta(y_t^* \mid x, y_{<t}^*)$\\[3pt]
          (Tại các vị trí $m_t = 1$)
        };

        \node[lossbox, text width=3.3cm, minimum height=1.5cm] (loss) at (13.6, 0) {
          \textbf{Hàm Mất mát SFT}\\
          $\mathcal{L}_{\text{SFT}}(\theta) = -\sum \log \pi_\theta$\\[3pt]
          \normalfont\scriptsize Chỉ tích lũy trên miền $y^*$
        };

        \draw[flow] (prompt.east) -- (llm.west |- prompt.east);
        \draw[flow] (resp.east) -- (llm.west |- resp.east);
        \draw[flow] (llm.east) -- (prob.west);
        \draw[flow] (prob.east) -- (loss.west);

        \coordinate (grad_start) at (13.6, -2.05);
        \coordinate (grad_end) at (4.7, -2.05);
        \draw[gradflow] (loss.south) -- (grad_start) -- node[below, font=\scriptsize\sffamily\bfseries, text=stanfordred] {Lan truyền ngược Gradient $\nabla_\theta \mathcal{L}_{\text{SFT}}$ cập nhật tham số $\theta$} (grad_end) -- (llm.south);
      \end{tikzpicture}%
    }
  \end{center}
  \vspace{-0.05cm}
  \footnotesize
  \textbf{Nguyên lý tính toán:} Truy vấn $x$ tham gia tính ma trận Key--Value Attention ở lượt truyền thuận, nhưng bị khóa đạo hàm ở lượt truyền ngược nhờ mặt nạ $m_t = 0$.
\end{frame}

% ----------------------------------------------------------------------------
% SLIDE: CÔNG THỨC TOÁN HỌC CỦA SFT & TẠI SAO CHỈ TÍNH LỖI TRÊN y*
% ----------------------------------------------------------------------------
\begin{frame}{Phần 1.2: Cơ sở Toán học của Tinh chỉnh có Giám sát (SFT)}{Ước lượng Hợp lý Cực đại có Điều kiện (Conditional Maximum Likelihood Estimation)}
  \begin{block}{Công thức 2: Hàm Mất mát Tinh chỉnh có Giám sát (Conditional MLE Loss)}
    \[
      \mathcal{L}_{\text{SFT}}(\theta) = -\E_{(x, y^*) \sim \mathcal{D}_{\text{SFT}}} \left[ \sum_{t=1}^{|y^*|} \log \pi_\theta\big(y_t^* \;\big|\; x, y_1^*, \dots, y_{t-1}^*\big) \right]
    \]
  \end{block}

  \begin{columns}[T,onlytextwidth]
    \begin{column}{0.51\textwidth}
      \subhead{Bản chất Giải tích}
      \small
      \begin{itemize}\setlength{\itemsep}{6pt}
        \item \textbf{Quy tắc dây chuyền xác suất:}
          $\log \pi_\theta(y^* \mid x) = \sum_{t=1}^{|y^*|} \log \pi_\theta(y_t^* \mid x, y_{<t}^*)$.
        \item \textbf{Tương đương Phân kỳ KL Thuận:}
          $\argmin_\theta \mathcal{L}_{\text{SFT}}(\theta) \equiv \argmin_\theta \DKL(P_{\text{data}} \parallel \pi_\theta)$.
      \end{itemize}
    \end{column}
    \begin{column}{0.46\textwidth}
      \limithead{Lý do loại trừ $\log \pi_\theta(x)$}
      \small
      \begin{itemize}\setlength{\itemsep}{6pt}
        \item Chỉ tối ưu hóa phân phối có điều kiện $P(Y \mid X)$, không mô hình hóa phân phối biên $P(X)$.
        \item Tập trung toàn bộ dung lượng tham số vào miền sinh phản hồi $y^*$.
      \end{itemize}
    \end{column}
  \end{columns}
\end{frame}

% ----------------------------------------------------------------------------
% SLIDE: GIẢ THUYẾT LIMA & LỌC MẪU REJECTION SAMPLING
% ----------------------------------------------------------------------------
\begin{frame}{Phần 1.2: Giả thuyết Căn chỉnh Bề mặt (LIMA) \& Lọc Mẫu Tổng hợp}{Vai trò của Chất lượng Dữ liệu (Zhou et al., 2023) và Hiện tượng Suy thoái Phân phối (Shumailov et al., 2024)}
  \begin{columns}[T,onlytextwidth]
    \begin{column}{0.48\textwidth}
      \motivehead{1. Giả thuyết LIMA (2023)}
      \footnotesize
      \begin{itemize}\setlength{\itemsep}{5pt}
        \item \textbf{Căn chỉnh bề mặt (Superficial Alignment):} Tri thức cốt lõi đã nằm trong Base Model; SFT chỉ ánh xạ biểu diễn nội tại ra định dạng giao tiếp chuẩn.
        \item \textbf{Ưu tiên chất lượng hơn số lượng:} $1{,}000$ mẫu chuẩn xác cao vượt trội so với $50{,}000$ mẫu thu thập đại trà chứa nhiễu.
      \end{itemize}
    \end{column}
    \begin{column}{0.48\textwidth}
      \limithead{2. Lọc Mẫu \& Suy thoái Phân phối}
      \footnotesize
      \begin{itemize}\setlength{\itemsep}{5pt}
        \item \textbf{Lọc mẫu theo ngưỡng (Rejection Sampling):} Sinh $N$ ứng viên $\{y_1, \dots, y_N\} \sim \pi_\theta$, dùng hàm đánh giá lọc lấy mẫu đạt điểm cao nhất đưa vào $\mathcal{D}_{\text{SFT}}$.
        \item \textbf{Suy thoái phân phối (Model Collapse):} Huấn luyện đệ quy trên dữ liệu tổng hợp không kiểm chứng làm co hẹp miền giá trị (\textit{Support Shrinkage}) và mất phương sai vùng đuôi.
      \end{itemize}
    \end{column}
  \end{columns}
\end{frame}

\subsection{1.3. Tổng kết Phần 1 (SFT): Thành tựu \& Các Vấn đề Tồn đọng cho Phần 2}

% ----------------------------------------------------------------------------
% SLIDE: SƠ ĐỒ TEACHER FORCING VS AUTOREGRESSIVE
% ----------------------------------------------------------------------------
\begin{frame}{Phần 1.3: Điểm Nghẽn Lý thuyết của SFT --- Sự Lệch Phân phối Tiếp xúc}{Đối chiếu Cơ chế Teacher Forcing lúc Huấn luyện và Giải mã Tự hồi quy lúc Suy diễn}
  \vspace{0.15cm}
  \begin{center}
    \resizebox{0.96\textwidth}{!}{%
      \begin{tikzpicture}
        \node[rewardbox, text width=3.2cm, minimum height=1.15cm] (tf_label) at (0, 1.05) {
          \textbf{1. Khi Huấn luyện}\\(Teacher Forcing)
        };
        \node[data, text width=2.3cm, minimum height=1.0cm] (tf1) at (3.5, 1.05) {Token chuẩn $y_1^*$};
        \node[data, text width=2.3cm, minimum height=1.0cm] (tf2) at (6.5, 1.05) {Token chuẩn $y_2^*$};
        \node[data, text width=2.3cm, minimum height=1.0cm] (tf3) at (9.5, 1.05) {Token chuẩn $y_3^*$};
        \node[rewardbox, text width=3.9cm, minimum height=1.15cm] (tf_res) at (13.5, 1.05) {
          Tiền tố điều kiện luôn thuộc\\phân phối chuẩn $y_{<t}^* \sim P_{\text{data}}$
        };

        \draw[flow] (tf_label.east) -- (tf1.west);
        \draw[flow] (tf1.east) -- (tf2.west);
        \draw[flow] (tf2.east) -- (tf3.west);
        \draw[flow] (tf3.east) -- (tf_res.west);

        \node[lossbox, text width=3.2cm, minimum height=1.15cm] (ar_label) at (0, -1.05) {
          \textbf{2. Khi Suy diễn}\\(Autoregressive Rollout)
        };
        \node[data, text width=2.3cm, minimum height=1.0cm] (ar1) at (3.5, -1.05) {Tự sinh $\hat{y}_1$};
        \node[lossbox, text width=2.3cm, minimum height=1.0cm] (ar2) at (6.5, -1.05) {Sai lệch $\hat{y}_2$};
        \node[lossbox, text width=2.3cm, minimum height=1.0cm] (ar3) at (9.5, -1.05) {Trượt quỹ đạo $\hat{y}_3$};
        \node[lossbox, text width=3.9cm, minimum height=1.15cm] (ar_res) at (13.5, -1.05) {
          Tích lũy sai số theo độ dài $T$:\\$P(\text{Quỹ đạo đúng}) = (1-\epsilon)^T \to 0$
        };

        \draw[flow] (ar_label.east) -- (ar1.west);
        \draw[gradflow] (ar1.east) -- (ar2.west);
        \draw[gradflow] (ar2.east) -- (ar3.west);
        \draw[gradflow] (ar3.east) -- (ar_res.west);
      \end{tikzpicture}%
    }
  \end{center}
  \vspace{0.08cm}
  \begin{alertblock}{Sự Dịch chuyển Phân phối Trạng thái (Covariate Shift -- Ross et al., 2011)}
    \footnotesize
    Khi huấn luyện: $s_t = (x, y_{<t}^*) \sim P_{\text{data}}$. Khi suy diễn: $\hat{s}_t = (x, \hat{y}_{<t}) \sim \pi_\theta$. Với sai số bước $\epsilon = 0.01$ và độ dài $T = 200$, xác suất giữ đúng quỹ đạo giảm còn $(0.99)^{200} \approx 13.4\%$.
  \end{alertblock}
\end{frame}

% ----------------------------------------------------------------------------
% SLIDE: TỔNG KẾT PHẦN 1 (1/2): NHỮNG GÌ ĐÃ ĐẠT ĐƯỢC & TỐI ƯU SO VỚI PRE-TRAINING
% ----------------------------------------------------------------------------
\begin{frame}{Phần 1.3: Tổng kết Phần 1 (1 / 2) --- Thành tựu Đạt được so với Tiền huấn luyện}{Đánh giá Đóng góp của Tinh chỉnh có Giám sát (2021--2022) so với Base Model (2020)}
  \motivehead{Những gì SFT Đã Giải quyết \& Tối ưu hơn so với Pre-training}
  \vspace{0.05cm}
  \footnotesize
  \renewcommand{\arraystretch}{1.32}
  \begin{tabularx}{\textwidth}{@{}p{3.4cm} X X@{}}
    \toprule
    \textbf{Tiêu chí Đánh giá} & \textbf{Tiền huấn luyện -- Base Model (2020)} & \textbf{Tinh chỉnh SFT (Phần 1 -- 2021--2022)} \\
    \midrule
    \textbf{1. Hàm mục tiêu} & $\mathcal{L}_{\text{MLE}}(u_t \mid u_{<t})$ trên toàn chuỗi & $\mathcal{L}_{\text{SFT}}(y_t^* \mid x, y_{<t}^*)$ với mặt nạ $m_t = 0$ trên $x$ \\
    \hline
    \textbf{2. Chế độ vận hành} & Chỉ nối tiếp văn bản, phụ thuộc mồi mẫu Few-shot & \textbf{\textcolor{tealgreen}{Khắc phục Objective Mismatch}}: Giải quyết chỉ thị trực tiếp ở chế độ Zero-shot \\
    \hline
    \textbf{3. Kiểm soát độ dài} & Sinh tràn lan tới giới hạn cửa sổ ngữ cảnh & Ước lượng chính xác xác suất dừng chuỗi tại \texttt{<EOS>} \\
    \hline
    \textbf{4. Độ phức tạp hệ thống} & Hàng nghìn tỷ token, chi phí tính toán khổng lồ & Chỉ cần $10^3\text{--}10^5$ mẫu trên \textbf{1 mô hình} ($\pi_\theta$), tối ưu hóa lồi cục bộ ổn định \\
    \bottomrule
  \end{tabularx}
\end{frame}

% ----------------------------------------------------------------------------
% SLIDE: TỔNG KẾT PHẦN 1 (2/2): VẤN ĐỀ TỒN ĐỌNG MÀ PHẦN 2 (RLHF + PPO) GIẢI QUYẾT
% ----------------------------------------------------------------------------
\begin{frame}{Phần 1.3: Tổng kết Phần 1 (2 / 2) --- Các Vấn đề Tồn đọng Dẫn tới Phần 2 (RLHF)}{Ba Giới hạn Cấu trúc của SFT mà Giai đoạn 2 (RLHF + PPO, 2022) Giải quyết}
  \limithead{Ba Bài toán Mở mà SFT Không thể Tự Giải quyết}
  \small
  \begin{enumerate}\setlength{\itemsep}{6pt}
    \item \textbf{Lệch Phân phối Tiếp xúc (Exposure Bias):}\\
      \footnotesize Mô hình chỉ quan sát tiền tố mẫu $y_{<t}^* \sim P_{\text{data}}$, chưa bao giờ được huấn luyện trên chuỗi tự sinh $\hat{y} \sim \pi_\theta$.
    \item \textbf{Kích thích Ảo giác Ngoại lai (Extrinsic Hallucination):}\\
      \footnotesize Cực đại hóa $\log \pi_\theta(y^* \mid x)$ trên cả tri thức nằm ngoài Base Model khiến mô hình suy giảm hiệu chỉnh xác suất (Miscalibration) và sinh thông tin bịa đặt.
    \item \textbf{Thiếu vắng Tín hiệu từ Mẫu Âm (Absence of Negative Samples):}\\
      \footnotesize SFT chỉ gia tăng xác suất cho mẫu dương chuẩn $y^*$ ($\uparrow$), không có đạo hàm phạt ($\downarrow$) đối với các mẫu âm sai lệch $y_l$.
  \end{enumerate}

  \vspace{0.06cm}
  \begin{exampleblock}{Bài toán Chuyển giao sang Phần 2 (RLHF \& PPO -- 2022)}
    \footnotesize
    Thiết kế cơ chế học từ \textbf{cặp đối chiếu gồm Mẫu dương $y_w$ và Mẫu âm $y_l$}, kết hợp lấy mẫu trực tuyến $\hat{y} \sim \pi_\theta$ để khử \textit{Exposure Bias}.
  \end{exampleblock}
\end{frame}

% ============================================================================
% PHẦN 2 (2022 - ĐẦU 2023): CĂN CHỈNH MÔ HÌNH BẰNG HỌC TĂNG CƯỜNG (RLHF & PPO)
% ============================================================================
\section{Phần 2 (2022--Đầu 2023): Căn chỉnh Mô hình bằng Học Tăng cường (RLHF \& PPO)}

\subsection{2.1. Dữ liệu Đối chiếu Cặp $(y_w, y_l)$ \& Mô hình Định giá Thưởng ($r_\phi$)}

% ----------------------------------------------------------------------------
% SLIDE: DỮ LIỆU ĐỐI CHIẾU CẶP: MẪU DƯƠNG y_w VÀ MẪU ÂM y_l
% ----------------------------------------------------------------------------
\begin{frame}{Phần 2.1: Dữ liệu Đối chiếu Cặp --- Mẫu Dương ($y_w$) và Mẫu Âm ($y_l$)}{Bước Ngoặt InstructGPT (Ouyang et al., 2022): Từ Học Mẫu Đơn sang Học Quan hệ Thứ tự}
  \subhead{Bất đối xứng Độ phức tạp giữa Bài toán Sinh (Generation) và Kiểm định (Verification)}
  \small
  Thay vì biên soạn nhãn chuẩn $y^*$ như SFT, mô hình sinh cặp ứng viên $(y_1, y_2) \sim \pi_{\text{SFT}}(\cdot \mid x)$ để xác lập quan hệ thứ tự $\mathbf{y_w \succ y_l}$ (\textbf{Mẫu dương $y_w$} ưu việt hơn \textbf{Mẫu âm $y_l$}):

  \vspace{0.08cm}
  \begin{columns}[T,onlytextwidth]
    \begin{column}{0.48\textwidth}
      \limithead{Hạn chế của Định lượng Vô hướng Tuyệt đối}
      \footnotesize
      \begin{itemize}\setlength{\itemsep}{6pt}
        \item \textbf{Phương sai cá nhân cao:} Thang điểm tuyệt đối ($1\text{--}10$) không đồng nhất giữa những người đánh giá.
        \item \textbf{Trôi chuẩn định lượng (Calibration Drift):} Ngưỡng chấm của cùng một chuyên gia biến thiên theo thời gian.
      \end{itemize}
    \end{column}
    \begin{column}{0.48\textwidth}
      \motivehead{Ưu thế của Cặp Đối chiếu $(y_w \succ y_l)$}
      \footnotesize
      \begin{itemize}\setlength{\itemsep}{6pt}
        \item \textbf{Bất biến với thang đo:} Quan hệ nhị phân $y_w \succ y_l$ có độ đồng thuận thống kê cao.
        \item \textbf{Bổ sung Mẫu âm $y_l$ (Negative Sample):} Cung cấp biên phân cách rõ ràng để triệt tiêu hành vi sai lệch.
      \end{itemize}
    \end{column}
  \end{columns}
\end{frame}

% ----------------------------------------------------------------------------
% SLIDE: SƠ ĐỒ KIẾN TRÚC & LUỒNG DỮ LIỆU HUẤN LUYỆN REWARD MODEL
% ----------------------------------------------------------------------------
\begin{frame}{Phần 2.1: Sơ đồ Kiến trúc \& Luồng Dữ liệu Huấn luyện Mô hình Thưởng ($r_\phi$)}{Học Hàm Định giá Vô hướng $r_\phi(x, y) \in \mathbb{R}$ từ Cặp Mẫu Dương $y_w$ và Mẫu Âm $y_l$}
  \vspace{0.15cm}
  \begin{center}
    \resizebox{0.96\textwidth}{!}{%
      \begin{tikzpicture}
        \node[data, text width=2.8cm, minimum height=1.2cm] (prompt) at (0, 0) {
          \textbf{Truy vấn $x$}\\
          (Từ tập $\mathcal{D}_{\text{pair}}$)
        };

        \node[data, fill=softgreen, draw=tealgreen, text width=3.3cm, minimum height=1.1cm] (yw) at (4.2, 1.15) {
          \textbf{Mẫu Dương $y_w$}\\
          (Positive Sample)
        };

        \node[data, fill=softred, draw=stanfordred, text width=3.3cm, minimum height=1.1cm] (yl) at (4.2, -1.15) {
          \textbf{Mẫu Âm $y_l$}\\
          (Negative Sample)
        };

        \node[trainmodel, text width=3.8cm, minimum height=3.1cm] (rm) at (9.0, 0) {
          \textbf{Mô hình Thưởng $r_\phi$}\\
          (Reward Model)\\[6pt]
          \normalfont\scriptsize Khởi tạo từ $\pi_{\text{SFT}}$;\\
          Thay đầu từ vựng $|\mathcal{V}|$ bằng\\
          \textbf{Đầu Chiếu Vô hướng}\\
          $W_r \in \mathbb{R}^{d \times 1}$ tại \texttt{<EOS>}
        };

        \node[rewardbox, text width=2.9cm, minimum height=1.0cm] (rw) at (13.3, 1.15) {
          Điểm mẫu dương\\$\mathbf{r_\phi(x, y_w) \in \mathbb{R}}$
        };

        \node[lossbox, fill=softred, draw=stanfordred, text width=2.9cm, minimum height=1.0cm] (rl) at (13.3, -1.15) {
          Điểm mẫu âm\\$\mathbf{r_\phi(x, y_l) \in \mathbb{R}}$
        };

        \node[lossbox, text width=3.6cm, minimum height=1.6cm] (btloss) at (17.5, 0) {
          \textbf{Hàm Mất mát Bradley-Terry}\\[3pt]
          $\mathcal{L}_{\text{RM}} = -\log \sigma(r_w - r_l)$\\[3pt]
          \normalfont\scriptsize Cực đại hóa hiệu số $r_w - r_l$
        };

        \coordinate (bus_p) at (1.9, 0);
        \draw[line width=1.05pt, draw=navyblue!85!black] (prompt.east) -- (bus_p);
        \draw[flow] (bus_p) |- (yw.west);
        \draw[flow] (bus_p) |- (yl.west);

        \draw[flow] (yw.east) -- ([yshift=1.15cm]rm.west);
        \draw[flow] (yl.east) -- ([yshift=-1.15cm]rm.west);

        \draw[flow] ([yshift=1.15cm]rm.east) -- (rw.west);
        \draw[flow] ([yshift=-1.15cm]rm.east) -- (rl.west);

        \draw[flow] (rw.east) -| (btloss.north);
        \draw[flow] (rl.east) -| (btloss.south);

        \draw[gradflow] (btloss.south) -- (17.5, -2.35) -- node[below, font=\scriptsize\sffamily\bfseries, text=stanfordred] {Lan truyền ngược Gradient $\nabla_\phi \mathcal{L}_{\text{RM}}$ cập nhật tham số $\phi$} (9.0, -2.35) -- (rm.south);
      \end{tikzpicture}%
    }
  \end{center}
  \vspace{-0.05cm}
  \footnotesize
  \textbf{Nguyên lý vận hành:} Cặp $(x, y_w)$ và $(x, y_l)$ đi qua cùng mạng $r_\phi$; đạo hàm đẩy điểm của mẫu dương $r_\phi(x, y_w)$ tăng lên ($\uparrow$) và điểm của mẫu âm $r_\phi(x, y_l)$ giảm xuống ($\downarrow$).
\end{frame}

% ----------------------------------------------------------------------------
% SLIDE: TOÁN HỌC CỦA REWARD MODELING — MÔ HÌNH BRADLEY-TERRY & BẤT BIẾN TỊNH TIẾN
% ----------------------------------------------------------------------------
\begin{frame}{Phần 2.1: Cơ sở Toán học của Mô hình Thưởng --- Mô hình Bradley-Terry}{Hàm Hợp lý Cực đại cho Quan hệ Thứ tự Cặp và Tính Bất biến Tịnh tiến}
  \begin{block}{Công thức 3: Mô hình Xác suất Bradley-Terry (1952) \& Hàm Mất mát $\mathcal{L}_{\text{RM}}(\phi)$}
    \footnotesize
    \[
      P(y_w \succ y_l \mid x) = \sigma\big(r_\phi(x, y_w) - r_\phi(x, y_l)\big), \qquad
      \mathcal{L}_{\text{RM}}(\phi) = -\E_{(x, y_w, y_l) \sim \mathcal{D}_{\text{pair}}} \Big[ \log \sigma\big(r_\phi(x, y_w) - r_\phi(x, y_l)\big) \Big]
    \]
  \end{block}

  \subhead{Tính chất Bất biến Tịnh tiến theo Trạng thái Truy vấn (Shift Invariance)}
  \small
  \begin{itemize}\setlength{\itemsep}{5pt}
    \item \textbf{Biểu diễn vô hướng tại \texttt{<EOS>}:} $r_\phi(x, y) = W_r^\top h_{\text{EOS}}(x, y) \in \mathbb{R}$ (do mặt nạ Causal Attention hội tụ toàn bộ tiền tố tại token cuối).
    \item \textbf{Bất biến tịnh tiến:} Với mọi hàm $C(x)$ chỉ phụ thuộc truy vấn $x$, phép biến đổi $\tilde{r}(x, y) = r(x, y) + C(x)$ bảo toàn hiệu số:
      \[
        \tilde{r}(x, y_w) - \tilde{r}(x, y_l) = r(x, y_w) - r(x, y_l) \implies P_{\tilde{r}}(y_w \succ y_l \mid x) \equiv P_r(y_w \succ y_l \mid x)
      \]
  \end{itemize}
\end{frame}

\subsection{2.2. Định luật Goodhart, Điều chuẩn Phân kỳ KL ($\beta \DKL$) \& Chu trình RL}

% ----------------------------------------------------------------------------
% SLIDE: ĐỊNH LUẬT GOODHART & REWARD HACKING -> ĐỘNG CƠ PHẠT KL
% ----------------------------------------------------------------------------
\begin{frame}{Phần 2.2: Định luật Goodhart \& Động cơ Điều chuẩn Phân kỳ Kullback-Leibler}{Tại sao Cực đại hóa Trực tiếp $\max_{\pi_\theta} \E[r_\phi(x, y)]$ Dẫn tới Suy biến Phân phối?}
  \begin{columns}[T,onlytextwidth]
    \begin{column}{0.48\textwidth}
      \limithead{Định luật Goodhart \& Trục lợi Phần thưởng}
      \footnotesize
      \begin{itemize}\setlength{\itemsep}{6pt}
        \item \textbf{Định luật Goodhart:} Khi hàm xấp xỉ $r_\phi$ bị tối ưu hóa cực đoan, nó mất tính tương quan với chất lượng thực tế.
        \item \textbf{Sai số ngoại suy (OOD Error):} $r_\phi$ chỉ chính xác trong lân cận phân phối của $\pi_{\text{SFT}}$.
        \item \textbf{Trục lợi phần thưởng (Reward Hacking -- Gao et al., 2023):} Chính sách $\pi_\theta$ dịch chuyển sang vùng ngoại lai để khai thác cực đại giả của $r_\phi$.
      \end{itemize}
    \end{column}
    \begin{column}{0.48\textwidth}
      \motivehead{Giải pháp: Điều chuẩn Vùng Tin cậy bằng $\DKL$}
      \footnotesize
      \begin{itemize}\setlength{\itemsep}{6pt}
        \item Cố định mô hình SFT làm \textbf{Mô hình Tham chiếu ($\pi_{\text{ref}}$)} và phạt độ lệch phân phối:
          \[
            \DKL\big(\pi_\theta(\cdot \mid x) \parallel \pi_{\text{ref}}(\cdot \mid x)\big) = \E_{y \sim \pi_\theta} \left[ \log \frac{\pi_\theta(y \mid x)}{\pi_{\text{ref}}(y \mid x)} \right]
          \]
        \item \textbf{Tác dụng:} Neo $\pi_\theta$ trong vùng tin cậy của $r_\phi$ và bảo toàn cấu trúc ngôn ngữ gốc.
      \end{itemize}
    \end{column}
  \end{columns}
\end{frame}

% ----------------------------------------------------------------------------
% SLIDE: CÔNG THỨC TOÁN HỌC CỦA RLHF CÓ RÀNG BUỘC PHẠT KL
% ----------------------------------------------------------------------------
\begin{frame}{Phần 2.2: Hàm Mục tiêu Tối ưu hóa Chính sách trong RLHF}{Bài toán Cực đại hóa Kỳ vọng Phần thưởng có Điều chuẩn Phân kỳ KL (Ouyang et al., 2022)}
  \begin{block}{Công thức 4: Hàm Mục tiêu Tối ưu hóa Chính sách RLHF (KL-Regularized RL Objective)}
    \[
      \max_{\pi_\theta} \; J_{\text{RLHF}}(\pi_\theta) = \E_{x \sim \mathcal{D},\; y \sim \pi_\theta(y \mid x)} \Big[ \underbrace{r_\phi(x, y)}_{\text{Điểm định giá } r_\phi} - \;\beta \underbrace{\log \frac{\pi_\theta(y \mid x)}{\pi_{\text{ref}}(y \mid x)}}_{\text{Điều chuẩn phân kỳ KL}} \Big]
    \]
  \end{block}

  \subhead{Tín hiệu Phần thưởng Hiệu dụng (Effective Reward) \& Hệ số Điều chuẩn $\beta > 0$}
  \small
  \begin{itemize}\setlength{\itemsep}{6pt}
    \item Phần thưởng tổng hợp tại mỗi bước lấy mẫu trực tuyến $y \sim \pi_\theta(\cdot \mid x)$:
      \[
        R_{\text{total}}(x, y) = r_\phi(x, y) - \beta \log \frac{\pi_\theta(y \mid x)}{\pi_{\text{ref}}(y \mid x)}
      \]
    \item \textbf{Giới hạn tiệm cận:} Khi $\beta \to 0 \implies$ suy biến do \textit{Reward Hacking}; khi $\beta \to +\infty \implies$ chính sách bất động ($\pi_\theta \equiv \pi_{\text{ref}}$).
  \end{itemize}
\end{frame}

% ----------------------------------------------------------------------------
% SLIDE: SƠ ĐỒ CHU TRÌNH HỌC TĂNG CƯỜNG TRỰC TUYẾN & KHẮC PHỤC EXPOSURE BIAS
% ----------------------------------------------------------------------------
\begin{frame}{Phần 2.2: Chu trình Học Tăng cường Trực tuyến (On-Policy RL Loop)}{Cơ chế Tự Sinh Mẫu $\hat{y} \sim \pi_\theta$ Khắc phục Triệt để Hiện tượng Exposure Bias của Phần 1}
  \vspace{0.12cm}
  \begin{center}
    \resizebox{0.96\textwidth}{!}{%
      \begin{tikzpicture}
        \node[data, text width=3.1cm, minimum height=1.3cm] (state) at (0, 0) {
          \textbf{Trạng thái $s_t$}\\
          Truy vấn $x$ \& tiền tố\\tự sinh $\hat{y}_{<t}$
        };

        \node[trainmodel, text width=3.5cm, minimum height=1.3cm] (policy) at (4.6, 0) {
          \textbf{Chính sách $\pi_\theta(\hat{y}_t \mid s_t)$}\\
          Mô hình Ngôn ngữ\\đang tối ưu hóa
        };

        \node[data, fill=softgreen, draw=tealgreen, text width=3.4cm, minimum height=1.3cm] (action) at (9.3, 0) {
          \textbf{Quỹ đạo Tự sinh}\\
          Rút mẫu $\hat{y} \sim \pi_\theta(\cdot \mid x)$\\
          (On-Policy Rollout)
        };

        \node[rewardbox, text width=3.5cm, minimum height=1.3cm] (env) at (14.0, 0) {
          \textbf{Định giá Phần thưởng}\\
          $R_{\text{total}} = r_\phi(x, \hat{y}) - \beta \DKL$
        };

        \node[lossbox, text width=7.8cm, minimum height=1.1cm] (update) at (9.3, -2.0) {
          \textbf{Hàm Lợi thế $A(x, \hat{y}) = R_{\text{total}}(x, \hat{y}) - b(x)$}\\[2pt]
          \normalfont\scriptsize Khi $A > 0$: Tăng xác suất $\pi_\theta(\hat{y} \mid x)$ ($\uparrow$) \quad $|$ \quad Khi $A < 0$: Giảm xác suất $\pi_\theta(\hat{y} \mid x)$ ($\downarrow$)
        };

        \draw[flow] (state.east) -- (policy.west);
        \draw[flow] (policy.east) -- (action.west);
        \draw[flow] (action.east) -- (env.west);

        \draw[flow] (env.south) |- (update.east);
        \draw[gradflow] (update.west) -| node[pos=0.25, below, font=\scriptsize\sffamily\bfseries, text=stanfordred] {Cập nhật Đạo hàm Chính sách $\nabla_\theta J(\theta)$} (policy.south);
      \end{tikzpicture}%
    }
  \end{center}
  \vspace{0.05cm}
  \footnotesize
  \textbf{Khắc phục Exposure Bias:} Chính sách $\pi_\theta$ được huấn luyện trực tiếp trên các trạng thái tự sinh $\hat{s}_t = (x, \hat{y}_{<t}) \sim \pi_\theta$ và nhận phạt âm ($A < 0$) ngay khi đi chệch quỹ đạo.
\end{frame}

\subsection{2.3. Đạo hàm Chính sách, Thuật toán PPO (Mạng Giá trị $V_\psi$) \& Hệ thống 4 Mô hình}

% ----------------------------------------------------------------------------
% SLIDE: ĐỊNH LÝ POLICY GRADIENT & ĐƯỜNG CƠ SỞ BASELINE
% ----------------------------------------------------------------------------
\begin{frame}{Phần 2.3: Định lý Đạo hàm Chính sách (Policy Gradient) \& Đường Cơ sở $b(x)$}{Giải quyết Tính Không Khả vi của Phép Rút mẫu Token Rời rạc (Williams, 1992)}
  \begin{block}{Công thức 5: Đạo hàm Chính sách (REINFORCE với Baseline) \& Hàm Lợi thế $A(x, y)$}
    \[
      \nabla_\theta J(\theta) = \E_{x \sim \mathcal{D},\; y \sim \pi_\theta(\cdot \mid x)} \Big[ \underbrace{\big(R_{\text{total}}(x, y) - b(x)\big)}_{\text{Hàm Lợi thế } A(x, y)} \, \nabla_\theta \log \pi_\theta(y \mid x) \Big]
    \]
  \end{block}

  \begin{columns}[T,onlytextwidth]
    \begin{column}{0.49\textwidth}
      \subhead{1. Đồng nhất thức Đạo hàm Logarit}
      \footnotesize
      \begin{itemize}\setlength{\itemsep}{4pt}
        \item Phép rút mẫu rời rạc $y \sim \pi_\theta$ không có đạo hàm $\frac{\partial y}{\partial \theta}$.
        \item Áp dụng $\nabla_\theta \pi_\theta = \pi_\theta \nabla_\theta \log \pi_\theta$:
          \[
            \nabla_\theta \sum_y \pi_\theta(y) R(y) = \E_{y \sim \pi_\theta}\big[R(y)\nabla_\theta \log \pi_\theta(y)\big]
          \]
      \end{itemize}
    \end{column}
    \begin{column}{0.48\textwidth}
      \motivehead{2. Vai trò của Đường cơ sở $b(x) = \E[R]$}
      \footnotesize
      \begin{itemize}\setlength{\itemsep}{4pt}
        \item Phân tách dấu giữa quỹ đạo trên trung bình ($A > 0, \uparrow$) và dưới trung bình ($A < 0, \downarrow$).
        \item Giảm phương sai ước lượng Monte Carlo mà không làm chệch kỳ vọng (do $\E[b(x)\nabla_\theta \log \pi_\theta] = 0$).
      \end{itemize}
    \end{column}
  \end{columns}
\end{frame}

% ----------------------------------------------------------------------------
% SLIDE: THUẬT TOÁN PPO: MẠNG CRITIC & CƠ CHẾ PPO-CLIP
% ----------------------------------------------------------------------------
\begin{frame}{Phần 2.3: Thuật toán Tối ưu hóa Chính sách Cận biên (PPO -- Schulman et al., 2017)}{Kết hợp Cắt tỉa Tỷ số Hợp lý ($\clip$) và Mạng Giá trị (Critic Network $V_\psi$)}
  \begin{block}{Công thức 6: Hàm Mục tiêu Cắt tỉa Cận biên (PPO-Clip Objective)}
    \[
      \mathcal{J}_{\text{PPO-Clip}}(\theta) = \E_{t} \Big[ \min \Big( \rho_t(\theta) \hat{A}_t, \; \clip\big(\rho_t(\theta), 1 - \epsilon, 1 + \epsilon\big) \hat{A}_t \Big) \Big],
      \qquad \rho_t(\theta) = \frac{\pi_\theta(y_t \mid x, y_{<t})}{\pi_{\theta_{\text{old}}}(y_t \mid x, y_{<t})}
    \]
  \end{block}

  \begin{columns}[T,onlytextwidth]
    \begin{column}{0.49\textwidth}
      \subhead{Cơ chế Chặn Biên $\clip(\rho_t, 1-\epsilon, 1+\epsilon)$}
      \footnotesize
      \begin{itemize}\setlength{\itemsep}{5pt}
        \item $\rho_t(\theta)$: Tỷ số xác suất của token $y_t$ giữa chính sách mới và chính sách lúc lấy mẫu.
        \item Khi $|\rho_t(\theta) - 1| > \epsilon$ (với $\epsilon = 0.2$), toán tử $\min(\cdot)$ triệt tiêu đạo hàm về $0$, ngăn bước cập nhật quá lớn.
      \end{itemize}
    \end{column}
    \begin{column}{0.47\textwidth}
      \subhead{Mạng Giá trị (Critic) $V_\psi(s_t)$}
      \footnotesize
      \begin{itemize}\setlength{\itemsep}{5pt}
        \item Ước lượng đường cơ sở kỳ vọng tại từng vị trí token $t$: $\hat{A}_t = R_t - V_\psi(s_t)$.
        \item Đòi hỏi duy trì một mạng nơ-ron $V_\psi$ có quy mô tham số tương đương mô hình chính $\pi_\theta$.
      \end{itemize}
    \end{column}
  \end{columns}
\end{frame}

% ----------------------------------------------------------------------------
% SLIDE: SƠ ĐỒ KIẾN TRÚC 4 MÔ HÌNH CỦA RLHF (PPO)
% ----------------------------------------------------------------------------
\begin{frame}{Phần 2.3: Sơ đồ Kiến trúc \& Luồng Dữ liệu của Hệ thống RLHF + PPO (4 Mô hình)}{Cấu trúc Phối hợp Bốn Mạng Nơ-ron Đồng thời trên Bộ nhớ GPU}
  \vspace{0.08cm}
  \begin{center}
    \resizebox{0.96\textwidth}{!}{%
      \begin{tikzpicture}
        \node[data, text width=2.4cm, minimum height=1.2cm] (prompt) at (0, 0) {
          \textbf{Truy vấn $x$}\\
          (Từ tập dữ liệu)
        };

        \node[trainmodel, text width=3.5cm, minimum height=1.35cm] (actor) at (4.2, 1.35) {
          \textbf{1. Actor / Policy $\pi_\theta$}\\
          (Đang huấn luyện)\\
          \normalfont\scriptsize Tự sinh chuỗi $y \sim \pi_\theta(\cdot \mid x)$
        };

        \node[frozenmodel, text width=3.5cm, minimum height=1.35cm] (ref) at (4.2, -1.35) {
          \textbf{2. Reference $\pi_{\text{ref}}$}\\
          (Cố định tham số)\\
          \normalfont\scriptsize Tính $\log \pi_{\text{ref}}(y \mid x)$
        };

        \node[data, fill=softgreen, draw=tealgreen, text width=2.5cm, minimum height=1.2cm] (resp) at (8.5, 0) {
          \textbf{Chuỗi tự sinh $y$}\\
          (On-policy Rollout)
        };

        \node[frozenmodel, text width=3.6cm, minimum height=1.35cm] (rm) at (12.8, 1.35) {
          \textbf{3. Reward Model $r_\phi$}\\
          (Cố định tham số)\\
          \normalfont\scriptsize Định giá $r_\phi(x, y)$
        };

        \node[trainmodel, fill=softyellow, draw=amberorange!90!black, text width=3.6cm, minimum height=1.35cm] (critic) at (12.8, -1.35) {
          \textbf{4. Critic / Value $V_\psi$}\\
          (Đang huấn luyện)\\
          \normalfont\scriptsize Ước lượng nền $V_\psi(s_t)$
        };

        \node[lossbox, text width=3.7cm, minimum height=2.2cm] (ppo) at (17.5, 0) {
          \textbf{Tối ưu hóa PPO}\\[3pt]
          $R_{\text{total}} = r_\phi - \beta \DKL$\\[2pt]
          $\hat{A}_t = R_{\text{total}} - V_\psi(s_t)$\\[3pt]
          \normalfont\scriptsize Cắt tỉa tỷ số $\clip(\rho_t)$
        };

        \coordinate (bus_p) at (1.8, 0);
        \draw[line width=1.05pt, draw=navyblue!85!black] (prompt.east) -- (bus_p);
        \draw[flow] (bus_p) |- (actor.west);
        \draw[flow] (bus_p) |- (ref.west);

        \draw[flow] (actor.east) -| (resp.north);
        \draw[flow] (resp.south) |- (ref.east);

        \coordinate (bus_r) at (10.3, 0);
        \draw[line width=1.05pt, draw=navyblue!85!black] (resp.east) -- (bus_r);
        \draw[flow] (bus_r) |- (rm.west);
        \draw[flow] (bus_r) |- (critic.west);

        \draw[flow] (rm.east) -| ([xshift=-0.4cm]ppo.north);
        \draw[flow] (critic.east) -| ([xshift=-0.4cm]ppo.south);
        \draw[flow] (ref.south) -- (4.2, -2.45) -- node[below, font=\scriptsize\sffamily, text=slate] {Truyền $\log \pi_{\text{ref}}(y \mid x)$ để tính phân kỳ $\beta \DKL(\pi_\theta \parallel \pi_{\text{ref}})$} (17.5, -2.45) -- (ppo.south);

        \draw[gradflow] (ppo.north) -- (17.5, 2.45) -- node[above, font=\scriptsize\sffamily\bfseries, text=stanfordred] {Cập nhật Gradient PPO đồng thời cho Policy $\pi_\theta$ và Critic $V_\psi$} (4.2, 2.45) -- (actor.north);
      \end{tikzpicture}%
    }
  \end{center}
\end{frame}

\subsection{2.4. Tổng kết Phần 2 (RLHF + PPO): Thành tựu \& Các Vấn đề Tồn đọng cho Phần 3}

% ----------------------------------------------------------------------------
% SLIDE: TỔNG KẾT PHẦN 2 (1/2): NHỮNG GÌ ĐÃ ĐẠT ĐƯỢC & TỐI ƯU SO VỚI SFT
% ----------------------------------------------------------------------------
\begin{frame}{Phần 2.4: Tổng kết Phần 2 (1 / 2) --- Thành tựu Đạt được so với SFT (Phần 1)}{Đánh giá Đóng góp của RLHF + PPO (2022) so với Tinh chỉnh có Giám sát (2021)}
  \motivehead{Những gì RLHF + PPO Đã Giải quyết \& Tối ưu hơn so với SFT}
  \vspace{0.05cm}
  \footnotesize
  \renewcommand{\arraystretch}{1.32}
  \begin{tabularx}{\textwidth}{@{}p{3.4cm} X X@{}}
    \toprule
    \textbf{Tiêu chí Đánh giá} & \textbf{Tinh chỉnh SFT (Phần 1 -- 2021)} & \textbf{RLHF + PPO (Phần 2 -- 2022)} \\
    \midrule
    \textbf{1. Phân phối trạng thái} & Học tĩnh trên tiền tố chuyên gia $y_{<t}^*$ (\textit{Teacher Forcing}) & \textbf{\textcolor{tealgreen}{Khử Exposure Bias}}: Học trực tuyến trên chuỗi tự sinh $\hat{y} \sim \pi_\theta$ (\textit{On-policy}) \\
    \hline
    \textbf{2. Tín hiệu đạo hàm} & Chỉ có tín hiệu dương ($\uparrow$) trên mẫu chuẩn $y^*$ & \textbf{\textcolor{tealgreen}{Đạo hàm hai chiều}}: Củng cố mẫu dương ($\hat{A}_t > 0, \uparrow$) và chế tài mẫu âm ($\hat{A}_t < 0, \downarrow$) \\
    \hline
    \textbf{3. Kiểm soát phân phối} & Suy giảm hiệu chỉnh xác suất khi ép học thuộc & Điều chuẩn vùng tin cậy $\beta\DKL(\pi_\theta \parallel \pi_{\text{ref}})$ chống suy biến \\
    \hline
    \textbf{4. Hiệu năng thực nghiệm} & Phụ thuộc mạnh vào kích thước mô hình nền & \textbf{InstructGPT 1.3B} vượt qua \textbf{GPT-3 175B} trên đánh giá mù \\
    \bottomrule
  \end{tabularx}
\end{frame}

% ----------------------------------------------------------------------------
% SLIDE: TỔNG KẾT PHẦN 2 (2/2): VẤN ĐỀ TỒN ĐỌNG MÀ PHẦN 3 (DPO) GIẢI QUYẾT
% ----------------------------------------------------------------------------
\begin{frame}{Phần 2.4: Tổng kết Phần 2 (2 / 2) --- Các Vấn đề Tồn đọng Dẫn tới Phần 3 (DPO)}{Ba Điểm Nghẽn Hệ thống \& Toán học của RLHF + PPO mà Giai đoạn 3 (DPO, 2023) Giải quyết}
  \limithead{Ba Rào cản Cốt lõi của Kiến trúc RLHF + PPO}
  \small
  \begin{enumerate}\setlength{\itemsep}{6pt}
    \item \textbf{Chi phí Bộ nhớ GPU Gấp 4 Lần (4-Model Memory Bottleneck):}\\
      \footnotesize Bắt buộc nạp đồng thời \textbf{4 mạng nơ-ron lớn} ($\pi_\theta, \pi_{\text{ref}}, r_\phi, V_\psi$) vào VRAM.
    \item \textbf{Nút thắt Thời gian Giải mã Tự hồi quy (Autoregressive Rollout Bottleneck):}\\
      \footnotesize Sinh mẫu tuần tự $\hat{y} \sim \pi_\theta$ tại mỗi bước cập nhật tiêu tốn $>70\%$ thời gian tính toán.
    \item \textbf{Phương sai Đạo hàm Cao \& Sai số Xấp xỉ Kép ($r_\phi, V_\psi$):}\\
      \footnotesize Sai số ước lượng từ cả mạng thưởng $r_\phi$ lẫn mạng giá trị $V_\psi$ gây mất ổn định hội tụ.
  \end{enumerate}

  \vspace{0.06cm}
  \begin{exampleblock}{Bài toán Chuyển giao sang Phần 3 (DPO -- 2023)}
    \footnotesize
    Giải bài toán tối ưu hóa RLHF bằng \textbf{nghiệm đóng giải tích} để cập nhật trực tiếp $\pi_\theta$ từ cặp mẫu dương/âm $(y_w, y_l)$ mà \textbf{bỏ hoàn toàn vòng lặp RL, bỏ cả $r_\phi$ lẫn $V_\psi$}.
  \end{exampleblock}
\end{frame}

% ============================================================================
% PHẦN 3 (2023 - 2024): TỐI ƯU HÓA CHÍNH SÁCH TRỰC TIẾP (DPO & CÁC BIẾN THỂ)
% ============================================================================
\section{Phần 3 (2023--2024): Tối ưu hóa Chính sách Trực tiếp (DPO \& Các Biến thể)}

\subsection{3.1. Từ Hàm Mục tiêu RLHF đến Hàm Mất mát DPO (Nghiệm Đóng Giải tích)}

% ----------------------------------------------------------------------------
% SLIDE: ĐỘNG CƠ CỦA DPO (05/2023)
% ----------------------------------------------------------------------------
\begin{frame}{Phần 3.1: Động cơ Toán học của Direct Preference Optimization (DPO)}{Thay thế Vòng lặp RL 4 Mô hình bằng Nghiệm Đóng Giải tích (Rafailov et al., 2023)}
  \subhead{Ý tưởng Đột phá: Ánh xạ Giải tích Trực tiếp giữa Hàm Thưởng $r(x, y)$ và Chính sách $\pi^*(y \mid x)$}
  \small
  \begin{itemize}\setlength{\itemsep}{5pt}
    \item \textbf{Quy trình gián tiếp của RLHF + PPO (Phần 2):}
      \[
        \text{Cặp } (y_w, y_l) \xrightarrow{\text{Bước 1 (MLE)}} \text{Hàm thưởng } r_\phi(x, y) \xrightarrow{\text{Bước 2 (PPO + } V_\psi\text{)}} \text{Chính sách } \pi_\theta(y \mid x)
      \]
    \item \textbf{Quy trình trực tiếp của DPO (Phần 3):} Giải nghiệm đóng của bài toán cực trị RLHF để biểu diễn ngược $r(x, y)$ theo $\pi_\theta(y \mid x)$, qua đó tối ưu hóa thẳng $\pi_\theta$ trên cặp mẫu dương/âm $(y_w, y_l)$!
  \end{itemize}

  \vspace{0.06cm}
  \footnotesize
  \renewcommand{\arraystretch}{1.28}
  \begin{tabularx}{\textwidth}{@{}l X X@{}}
    \toprule
    \textbf{Tiêu chí Hệ thống} & \textbf{RLHF + PPO (Phần 2 -- 2022)} & \textbf{DPO (Phần 3 -- 2023)} \\
    \midrule
    \textbf{Số mô hình trong VRAM} & \textbf{4 mạng:} $\pi_\theta, \pi_{\text{ref}}, r_\phi, V_\psi$ & \textbf{2 mạng:} Chỉ cần $\pi_\theta$ và $\pi_{\text{ref}}$ \\
    \hline
    \textbf{Giải mã tự hồi quy lúc học} & Bắt buộc ở mọi bước cập nhật (\textit{Rollout}) & \textbf{Không cần} (Tính log-prob song song) \\
    \bottomrule
  \end{tabularx}
\end{frame}

% ----------------------------------------------------------------------------
% SLIDE: TỪ MỤC TIÊU RLHF (B1) ĐẾN HÀM MẤT MÁT DPO (B7)
% ----------------------------------------------------------------------------
\begin{frame}{Phần 3.1: Chuyển đổi Giải tích từ Mục tiêu RLHF sang Hàm Mất mát DPO}{Từ Bài toán Cực trị có Ràng buộc KL (B1) đến Hàm Mất mát Trực tiếp (B7)}
  \begin{columns}[T,onlytextwidth]
    \begin{column}{0.48\textwidth}
      \subhead{1. Mục tiêu RLHF \& Nghiệm Đóng (B1)}
      \footnotesize
      \textbf{Bài toán cực trị có điều chuẩn KL:}
      \[
        \max_{\pi} \E_{y \sim \pi} \left[ r(x, y) - \beta \log \frac{\pi(y \mid x)}{\pi_{\text{ref}}(y \mid x)} \right]
      \]
      \textbf{Nghiệm cực trị Gibbs--Boltzmann (BĐT Gibbs):}
      \[
        \pi^*(y \mid x) = \frac{1}{Z(x)} \pi_{\text{ref}}(y \mid x) \exp\left(\frac{1}{\beta} r(x, y)\right)
      \]
      \textbf{Biểu diễn ngược hàm thưởng:}
      \[
        \boxed{r(x, y) = \beta \log \frac{\pi^*(y \mid x)}{\pi_{\text{ref}}(y \mid x)} + \beta \log Z(x)}
      \]
      \textit{\scriptsize($Z(x) = \sum_{y} \pi_{\text{ref}}\exp(\frac{1}{\beta}r)$: hàm phân hoạch)}
    \end{column}

    \begin{column}{0.49\textwidth}
      \subhead{2. Khử $Z(x)$ \& Hàm Mất mát DPO (B7)}
      \footnotesize
      \textbf{Thế vào xác suất cặp Bradley--Terry:}\\
      Khi lấy hiệu số $r(x, y_w) - r(x, y_l)$, đại lượng $\beta \log Z(x)$ \textbf{tự triệt tiêu hoàn toàn}:
      \[
        r(x, y_w) - r(x, y_l) = \beta \log \frac{\pi^*(y_w \mid x)}{\pi_{\text{ref}}(y_w \mid x)} - \beta \log \frac{\pi^*(y_l \mid x)}{\pi_{\text{ref}}(y_l \mid x)}
      \]
      \vspace{-0.15cm}
      \begin{block}{Công thức 7: Hàm Mất mát DPO (B7)}
        \scriptsize
        Thay $\pi^*$ bằng $\pi_\theta$, tối ưu MLE trên tập cặp:
        \[
          \mathcal{L}_{\text{DPO}}(\theta) = -\E_{(x, y_w, y_l)} \left[ \log \sigma \left( \hat{r}_\theta(x, y_w) - \hat{r}_\theta(x, y_l) \right) \right]
        \]
        với phần thưởng ngầm định:
        $\hat{r}_\theta(x, y) = \beta \log \frac{\pi_\theta(y \mid x)}{\pi_{\text{ref}}(y \mid x)}$.
      \end{block}
    \end{column}
  \end{columns}
\end{frame}

\subsection{3.2. Cơ chế Gradient, Kiến trúc 2 Mô hình \& Các Biến thể (2024)}

% ----------------------------------------------------------------------------
% SLIDE: CƠ CHẾ GRADIENT VÀ TỰ ĐIỀU TRỌNG SỐ CỦA DPO
% ----------------------------------------------------------------------------
\begin{frame}{Phần 3.2: Cơ chế Cập nhật Gradient và Tự Điều trọng số của DPO}{Cơ chế Tự Điều chỉnh Biên độ Cập nhật giữa Mẫu Dương ($y_w$) và Mẫu Âm ($y_l$)}
  \subhead{Biểu thức Gradient Đầy đủ của Hàm Mất mát DPO (Rafailov et al., 2023)}
  \small
  \[
    \nabla_\theta \mathcal{L}_{\text{DPO}}(\theta) = -\beta \E_{(x, y_w, y_l)} \Big[ \underbrace{\sigma\big(\hat{r}_\theta(x, y_l) - \hat{r}_\theta(x, y_w)\big)}_{\text{Hệ số lỗi xếp hạng } \in (0, 1)} \cdot \Big( \underbrace{\nabla_\theta \log \pi_\theta(y_w \mid x)}_{\text{Kéo Mẫu dương } y_w \; (\uparrow)} - \underbrace{\nabla_\theta \log \pi_\theta(y_l \mid x)}_{\text{Đẩy Mẫu âm } y_l \; (\downarrow)} \Big) \Big]
  \]

  \subhead{Ý nghĩa Cơ chế Tự Điều trọng số Động (Dynamic Weighting)}
  \footnotesize
  \begin{itemize}\setlength{\itemsep}{4pt}
    \item \textbf{Khi mô hình xếp hạng sai ($\hat{r}_\theta(x, y_l) > \hat{r}_\theta(x, y_w)$):}\\
      Hệ số lỗi $\sigma(\hat{r}_l - \hat{r}_w) > 0.5$ (tiệm cận $1.0$), tự động \textbf{khuếch đại biên độ gradient} để ép mạng sửa sai triệt để.
    \item \textbf{Khi mô hình đã phân loại đúng ($\hat{r}_\theta(x, y_w) \gg \hat{r}_\theta(x, y_l)$):}\\
      Hệ số $\sigma(\hat{r}_l - \hat{r}_w) \to 0$, gradient tự triệt tiêu $\to$ \textbf{ngăn ngừa quá khớp} trên các mẫu đã học tốt.
    \item \textbf{Tác động đối ngẫu:}\\
      Đồng thời kéo tăng hợp lý của mẫu dương $y_w$ và triệt hạ xác suất phát sinh mẫu âm $y_l$.
  \end{itemize}
\end{frame}

% ----------------------------------------------------------------------------
% SLIDE: SƠ ĐỒ KIẾN TRÚC & LUỒNG DỮ LIỆU CỦA DPO
% ----------------------------------------------------------------------------
\begin{frame}{Phần 3.2: Sơ đồ Kiến trúc \& Luồng Dữ liệu của Thuật toán DPO}{Tinh giản Hệ thống từ Bốn Mô hình (RLHF + PPO) xuống Hai Mô hình ($\pi_\theta, \pi_{\text{ref}}$)}
  \vspace{0.15cm}
  \begin{center}
    \resizebox{0.96\textwidth}{!}{%
      \begin{tikzpicture}
        \node[data, text width=3.3cm, minimum height=1.3cm] (data) at (0, 0) {
          \textbf{Cặp Đối chiếu Tĩnh}\\
          $(x, y_w, y_l) \sim \mathcal{D}_{\text{pair}}$\\
          \normalfont\tiny ($y_w$: Mẫu dương, $y_l$: Mẫu âm)
        };

        \node[trainmodel, text width=3.8cm, minimum height=1.35cm] (actor) at (5.0, 1.2) {
          \textbf{1. Policy Model $\pi_\theta$}\\
          (Đang tối ưu hóa)\\
          \normalfont\scriptsize Tính $\log \pi_\theta(y_w \mid x), \;\log \pi_\theta(y_l \mid x)$
        };

        \node[frozenmodel, text width=3.8cm, minimum height=1.35cm] (ref) at (5.0, -1.2) {
          \textbf{2. Reference Model $\pi_{\text{ref}}$}\\
          (Cố định tham số)\\
          \normalfont\scriptsize Tính $\log \pi_{\text{ref}}(y_w \mid x), \;\log \pi_{\text{ref}}(y_l \mid x)$
        };

        \node[rewardbox, text width=3.7cm, minimum height=2.1cm] (implicit) at (10.1, 0) {
          \textbf{Thưởng Ngầm định}\\[4pt]
          $\hat{r}_w = \beta \log \dfrac{\pi_\theta(y_w \mid x)}{\pi_{\text{ref}}(y_w \mid x)}$\\[4pt]
          $\hat{r}_l = \beta \log \dfrac{\pi_\theta(y_l \mid x)}{\pi_{\text{ref}}(y_l \mid x)}$
        };

        \node[lossbox, text width=3.5cm, minimum height=2.1cm] (dpoloss) at (14.8, 0) {
          \textbf{Hàm Mất mát DPO}\\[4pt]
          $\mathcal{L}_{\text{DPO}} = -\log \sigma(\hat{r}_w - \hat{r}_l)$\\[4pt]
          \normalfont\scriptsize Khử hoàn toàn $r_\phi$\\
          và mạng giá trị $V_\psi$
        };

        \coordinate (bus_in) at (2.25, 0);
        \draw[line width=1.05pt, draw=navyblue!85!black] (data.east) -- (bus_in);
        \draw[flow] (bus_in) |- (actor.west);
        \draw[flow] (bus_in) |- (ref.west);

        \draw[flow] (actor.east) -- (7.55, 1.2) |- ([yshift=0.45cm]implicit.west);
        \draw[flow] (ref.east) -- (7.55, -1.2) |- ([yshift=-0.45cm]implicit.west);
        \draw[flow] (implicit.east) -- (dpoloss.west);

        \draw[gradflow] (dpoloss.north) -- (14.8, 2.45) -- node[above, font=\scriptsize\sffamily\bfseries, text=stanfordred] {Lan truyền ngược Gradient trực tiếp cập nhật $\pi_\theta$ (không cần Rollout)} (5.0, 2.45) -- (actor.north);
      \end{tikzpicture}%
    }
  \end{center}
  \vspace{0.05cm}
  \footnotesize
  \textbf{Ưu thế hệ thống:} Tính toán log-xác suất song song trên toàn chuỗi $(y_w, y_l)$ trong một lượt truyền thuận duy nhất.
\end{frame}

% ----------------------------------------------------------------------------
% SLIDE: CÁC BIẾN THỂ HẬU DPO (2024: SIMPO, RPO, IPO, KTO)
% ----------------------------------------------------------------------------
\begin{frame}{Phần 3.2: Các Biến thể Mở rộng của DPO (Giai đoạn 2024)}{Khắc phục Thiên kiến Độ dài và Hiện tượng Dịch chuyển Hợp lý (SimPO, RPO, IPO, KTO)}
  \begin{columns}[T,onlytextwidth]
    \begin{column}{0.48\textwidth}
      \limithead{Ba Hạn chế Giải tích của DPO Chuẩn}
      \footnotesize
      \begin{enumerate}\setlength{\itemsep}{5pt}
        \item \textbf{Thiên kiến độ dài (Length Bias):}\\
          Tổng $\log \pi_\theta(y \mid x) = \sum_{t=1}^{|y|} \log \pi_\theta(y_t \mid x, y_{<t})$ tỷ lệ thuận với $|y|$, gây ưu tiên mẫu dài dòng.
        \item \textbf{Dịch chuyển hợp lý (Likelihood Displacement):}\\
          Chỉ tối ưu hiệu số $\hat{r}_w - \hat{r}_l$, có thể làm cả $\log \pi_\theta(y_w)$ lẫn $\log \pi_\theta(y_l)$ cùng giảm.
        \item \textbf{Đòi hỏi cặp đồng truy vấn $(x, y_w, y_l)$:}\\
          Không tận dụng được dữ liệu đánh giá nhị phân đơn lẻ.
      \end{enumerate}
    \end{column}
    \begin{column}{0.48\textwidth}
      \motivehead{Các Phương pháp Cải tiến (2024)}
      \footnotesize
      \begin{itemize}\setlength{\itemsep}{5pt}
        \item \textbf{SimPO (Meng et al., 2024):}\\
          Chuẩn hóa độ dài $\frac{\beta}{|y|}\log \pi_\theta(y \mid x)$ và thay $\pi_{\text{ref}}$ bằng hằng số biên $\gamma > 0$ (chỉ cần \textbf{1 mô hình}).
        \item \textbf{RPO / CPO (2024):}\\
          Cộng thêm $\alpha \mathcal{L}_{\text{SFT}}(y_w)$ để neo xác suất tuyệt đối của mẫu dương $y_w$.
        \item \textbf{IPO (Azar et al., 2024):}\\
          Hàm mất mát bình phương ngăn quá khớp khi xác suất cặp tiệm cận $1$.
        \item \textbf{KTO (Ethayarajh et al., 2024):}\\
          Học trực tiếp từ mẫu dương hoặc mẫu âm đơn lẻ.
      \end{itemize}
    \end{column}
  \end{columns}
\end{frame}

\subsection{3.3. Tổng kết Phần 3 (DPO): Thành tựu \& Các Vấn đề Tồn đọng cho Phần 4}

% ----------------------------------------------------------------------------
% SLIDE: TỔNG KẾT PHẦN 3 (1/2): NHỮNG GÌ ĐÃ ĐẠT ĐƯỢC & TỐI ƯU SO VỚI RLHF+PPO
% ----------------------------------------------------------------------------
\begin{frame}{Phần 3.3: Tổng kết Phần 3 (1 / 2) --- Thành tựu Đạt được so với RLHF + PPO}{Đánh giá Bước Tiến Toán học của DPO (2023--2024) so với Hệ thống PPO (2022)}
  \motivehead{Những gì DPO Đã Giải quyết \& Tối ưu hơn so với RLHF + PPO (Phần 2)}
  \vspace{0.05cm}
  \footnotesize
  \renewcommand{\arraystretch}{1.32}
  \begin{tabularx}{\textwidth}{@{}p{3.4cm} X X@{}}
    \toprule
    \textbf{Tiêu chí Đánh giá} & \textbf{RLHF + PPO (Phần 2 -- 2022)} & \textbf{DPO \& Biến thể (Phần 3 -- 2023--2024)} \\
    \midrule
    \textbf{1. Cơ sở giải tích} & Xấp xỉ lặp qua đạo hàm Monte Carlo và mạng $r_\phi, V_\psi$ & \textbf{\textcolor{tealgreen}{Nghiệm đóng Boltzmann}}: Khử hoàn toàn hằng số phân hoạch $\beta \log Z(x)$ \\
    \hline
    \textbf{2. Bộ nhớ GPU} & \textbf{4 mạng:} $\pi_\theta, \pi_{\text{ref}}, r_\phi, V_\psi$ & \textbf{\textcolor{tealgreen}{2 mạng}} ($\pi_\theta, \pi_{\text{ref}}$) với DPO, hoặc \textbf{1 mạng} với SimPO \\
    \hline
    \textbf{3. Tốc độ tính toán} & Chậm do giải mã tự hồi quy tuần tự (\textit{Rollout}) & \textbf{\textcolor{tealgreen}{Tăng tốc vượt bậc}}: Truyền thuận song song trên cặp $(y_w, y_l)$ \\
    \hline
    \textbf{4. Tính ổn định} & Phương sai đạo hàm cao, nhạy cảm siêu tham số & Hội tụ ổn định tương đương SFT (ứng dụng trên Llama 3, Qwen 2, Zephyr) \\
    \bottomrule
  \end{tabularx}
\end{frame}

% ----------------------------------------------------------------------------
% SLIDE: TỔNG KẾT PHẦN 3 (2/2): VẤN ĐỀ TỒN ĐỌNG MÀ PHẦN 4 (RLVR + GRPO) GIẢI QUYẾT
% ----------------------------------------------------------------------------
\begin{frame}{Phần 3.3: Tổng kết Phần 3 (2 / 2) --- Các Vấn đề Tồn đọng Dẫn tới Phần 4 (RLVR + GRPO)}{Ba Giới hạn của DPO \& Định giá Thần kinh mà Giai đoạn 4 (2024--2026) Giải quyết}
  \limithead{Tại sao Cuối Năm 2024 Giới Nghiên cứu Phải Quay lại Online RL cho Bài toán Suy luận?}
  \small
  \begin{enumerate}\setlength{\itemsep}{6pt}
    \item \textbf{Giới hạn Học Ngoại tuyến (Off-Policy Exploration Bottleneck):}\\
      \footnotesize DPO chỉ tối ưu trên tập cặp tĩnh $\mathcal{D}_{\text{pair}}$ thu thập trước; không thể tự khám phá các quỹ đạo suy luận mới nằm ngoài tập dữ liệu.
    \item \textbf{Thất bại của Định giá Thần kinh trên Toán học \& Thuật toán:}\\
      \footnotesize Trong bài toán suy luận đa bước, cả DPO lẫn mạng $r_\phi$ đều bị đánh lừa bởi các lời giải trình bày đẹp nhưng sai kết quả (\textit{Reward Hacking}).
    \item \textbf{Điểm nghẽn Mạng Giá trị $V_\psi$ khi Quay lại Online RL:}\\
      \footnotesize Nếu dùng lại PPO (Phần 2) cho chuỗi suy luận dài $10{,}000$ token, mạng Critic $V_\psi$ bị suy biến do phần thưởng thưa ở cuối chuỗi.
  \end{enumerate}

  \vspace{0.06cm}
  \begin{exampleblock}{Bài toán Chuyển giao sang Phần 4 (RLVR + GRPO -- Cuối 2024--2026)}
    \footnotesize
    Kết hợp \textbf{Bộ kiểm chứng tất định (RLVR)} chống \textit{Reward Hacking} và giải thuật \textbf{GRPO} chuẩn hóa nhóm $G$ quỹ đạo để \textbf{xóa bỏ mạng Critic $V_\psi$}.
  \end{exampleblock}
\end{frame}

% ============================================================================
% PHẦN 4 (CUỐI 2024 - 2026): HỌC TĂNG CƯỜNG KIỂM CHỨNG & SUY LUẬN (RLVR & GRPO)
% ============================================================================
\section{Phần 4 (Cuối 2024--2026): Học Tăng cường Kiểm chứng \& Suy luận (RLVR \& GRPO)}

\subsection{4.1. Phần thưởng Kiểm chứng Tất định (RLVR) \& Bài toán Gán Tín dụng (ORM vs.\ PRM)}

% ----------------------------------------------------------------------------
% SLIDE: TỪ NEURAL RM SANG RLVR (CUỐI 2024 - 2025)
% ----------------------------------------------------------------------------
\begin{frame}{Phần 4.1: Bước Ngoặt Cuối Năm 2024 -- 2025 --- Từ $r_\phi$ sang Kiểm chứng Tất định (RLVR)}{Khắc phục Sai số Xấp xỉ của Mạng Phần thưởng Thần kinh trong Toán học \& Thuật toán}
  \begin{columns}[T,onlytextwidth]
    \begin{column}{0.48\textwidth}
      \limithead{Hạn chế của Mạng Thưởng $r_\phi$ trong Suy luận}
      \footnotesize
      \begin{itemize}\setlength{\itemsep}{6pt}
        \item \textbf{Thiên kiến hình thức bề mặt:} Mạng $r_\phi$ gán điểm cao cho lời giải trình bày LaTeX chặt chẽ dù biến đổi sai ở bước cuối.
        \item \textbf{Trục lợi phần thưởng trên chuỗi dài:} Khi chạy RL nhiều bước, $\pi_\theta$ sinh lập luận ngụy biện để đánh lừa $r_\phi$ thay vì giải đúng bài toán.
      \end{itemize}
    \end{column}
    \begin{column}{0.48\textwidth}
      \motivehead{Cơ chế RLVR (RL with Verifiable Rewards)}
      \footnotesize
      \begin{itemize}\setlength{\itemsep}{6pt}
        \item Trong miền hình thức (Toán học, Lập trình, Logic), \textbf{tính đúng đắn của nghiệm là khách quan và tất định}.
        \item Thay thế mạng $r_\phi$ bằng \textbf{Bộ Kiểm chứng Tất định}:
          \begin{itemize}\scriptsize\setlength{\itemsep}{3pt}
            \item \textit{Toán học:} Đối sánh đại số biểu thức trong \texttt{\textbackslash boxed\{\}} với $y_{\text{gold}}$.
            \item \textit{Lập trình:} Biên dịch và thực thi bộ kiểm thử (Unit Tests).
          \end{itemize}
      \end{itemize}
    \end{column}
  \end{columns}
\end{frame}

% ----------------------------------------------------------------------------
% SLIDE: CƠ CHẾ HOẠT ĐỘNG CỦA RLVR
% ----------------------------------------------------------------------------
\begin{frame}{Phần 4.1: Hàm Phần thưởng Kiểm chứng Tất định (RLVR)}{Triệt tiêu Hiện tượng Trục lợi Phần thưởng Văn phong (Lambert et al., 2024; DeepSeek-AI, 2025)}
  \begin{block}{Hàm Phần thưởng Nhị phân Tất định (Deterministic Verifiable Reward)}
    \[
      R_{\text{RLVR}}(x, y) = \begin{cases}
        1 & \text{nếu } \mathrm{Verify}\big(\mathrm{ExtractAnswer}(y),\; y_{\text{gold}}\big) = \mathrm{True} \\
        0 & \text{trong mọi trường hợp còn lại}
      \end{cases}
    \]
  \end{block}

  \subhead{Đặc tính Giải tích của RLVR trong Tối ưu hóa Dài hạn}
  \small
  \begin{itemize}\setlength{\itemsep}{6pt}
    \item \textbf{Bất biến với nhiễu bề mặt:} Bất kể độ dài hay văn phong lập luận, sai lệch kết quả cuối cùng hoặc trượt kiểm thử đều nhận $R_{\text{RLVR}} = 0$.
    \item \textbf{Không có sai số ngoại suy (Zero OOD Degradation):} Cho phép duy trì vòng lặp RL qua hàng chục nghìn bước cập nhật để khám phá chuỗi suy luận sâu (như OpenAI o1/o3, DeepSeek-R1).
  \end{itemize}
\end{frame}

% ----------------------------------------------------------------------------
% SLIDE: ĐỘ MỊN GÁN TÍN DỤNG: ORM VS PRM (P1)
% ----------------------------------------------------------------------------
\begin{frame}{Phần 4.1: Bài toán Gán Tín dụng (Credit Assignment) --- ORM so với PRM}{Đối chiếu Giám sát Cấp độ Kết quả (Outcome) và Giám sát Cấp độ Quá trình (Process -- Lightman et al., 2023)}
  \subhead{Định vị Sai số trên Quỹ đạo Suy luận Đa bước $y = (s_1, s_2, \dots, s_K)$}
  \vspace{0.05cm}
  \begin{columns}[T,onlytextwidth]
    \begin{column}{0.48\textwidth}
      \limithead{1. Mô hình Thưởng Kết quả (ORM)}
      \footnotesize
      \begin{itemize}\setlength{\itemsep}{5pt}
        \item \textbf{Cơ chế:} Chỉ cấp một điểm thưởng duy nhất tại cuối chuỗi $y$.
        \item \textbf{Ưu điểm:} Không cần nhãn bước trung gian; tương thích trực tiếp với bộ kiểm chứng RLVR.
        \item \textbf{Hạn chế lý thuyết:} Tín hiệu thưa (\textit{Sparse Reward}) --- phạt đồng loạt mọi bước trong quỹ đạo sai.
      \end{itemize}
    \end{column}
    \begin{column}{0.48\textwidth}
      \motivehead{2. Mô hình Thưởng Quá trình (PRM)}
      \footnotesize
      \begin{itemize}\setlength{\itemsep}{5pt}
        \item \textbf{Cơ chế:} Cấp điểm thưởng $r_k$ ngay sau từng bước biến đổi trung gian $s_k$.
        \item \textbf{Ưu điểm lý thuyết:} Tín hiệu dày đặc (\textit{Dense Credit Assignment}) --- bảo toàn các bước đầu đúng và định vị chính xác bước phát sinh lỗi.
      \end{itemize}
    \end{column}
  \end{columns}
\end{frame}

% ----------------------------------------------------------------------------
% SLIDE: TẠI SAO PRM HẠN CHẾ TRÊN QUY MÔ LỚN -> CHỌN RULE-BASED ORM (P2)
% ----------------------------------------------------------------------------
\begin{frame}{Phần 4.1: Giới hạn Thực nghiệm của PRM và Sự Dịch chuyển về ORM dựa trên Luật}{Lý do các Hệ thống Suy luận Quy mô Lớn (DeepSeek-R1, 2025) Loại bỏ Mạng PRM}
  \limithead{Ba Rào cản Cấu trúc khiến Mạng PRM Suy thoái trong Huấn luyện RL Dài hạn}
  \small
  \begin{enumerate}\setlength{\itemsep}{7pt}
    \item \textbf{Bất định ranh giới bước (Step Boundary Ambiguity):}\\
      \footnotesize Tư duy tự do (thử nghiệm, tự phản biện, quay lui) không thể phân tách cứng nhắc thành các bước đồng nhất.
    \item \textbf{Chi phí gán nhãn quá trình đắt đỏ:}\\
      \footnotesize Bộ kiểm chứng tất định không chấm được bước nháp dở dang; huấn luyện mạng PRM đòi hỏi gán nhãn thủ công hoặc chạy tìm kiếm cây Monte Carlo (MCTS) tốn kém.
    \item \textbf{Trục lợi phần thưởng cấp độ bước (Step-level Reward Hacking):}\\
      \footnotesize Do PRM vẫn là mạng xấp xỉ thần kinh, chính sách $\pi_\theta$ học cách chia lời giải thành hàng trăm bước biến đổi hiển nhiên (\textit{``$1=1 \implies 2=2$''}) để tích lũy điểm dương giả tạo.
  \end{enumerate}
\end{frame}

\subsection{4.2. Thuật toán Tối ưu hóa Tương đối Nhóm (GRPO) \& Các Tinh chỉnh Mới}

% ----------------------------------------------------------------------------
% SLIDE: GIỚI HẠN CỦA PPO TRONG SUY LUẬN DÀI -> ĐỘNG CƠ RA ĐỜI GRPO
% ----------------------------------------------------------------------------
\begin{frame}{Phần 4.2: Giới hạn của PPO (Phần 2) trên Chuỗi Dài $\longrightarrow$ Động cơ của GRPO}{Từ Ước lượng Hàm Giá trị Tham số $V_\psi(s_t)$ sang Chuẩn hóa Thống kê Tương đối theo Nhóm}
  \begin{columns}[T,onlytextwidth]
    \begin{column}{0.48\textwidth}
      \limithead{Hai Giới hạn của PPO trên Chuỗi Suy luận Dài}
      \footnotesize
      \begin{enumerate}\setlength{\itemsep}{6pt}
        \item \textbf{Chi phí bộ nhớ gấp đôi từ mạng Critic $V_\psi$:}\\
          Phải lưu trữ và cập nhật mạng $V_\psi$ (70B--671B tham số) song song với $\pi_\theta$.
        \item \textbf{Phương sai cao của $V_\psi(s_t)$ dưới phần thưởng thưa:}\\
          Trong RLVR, điểm thưởng $\{0, 1\}$ chỉ xuất hiện ở token cuối thứ $10{,}000$; việc ước lượng $V_\psi(s_t)$ tại token $t \ll T$ có sai số rất lớn, làm lệch lợi thế $\hat{A}_t$.
      \end{enumerate}
    \end{column}
    \begin{column}{0.48\textwidth}
      \motivehead{Giải pháp GRPO (Shao et al., 2024 -- DeepSeekMath)}
      \footnotesize
      \textbf{Nguyên lý: Khử hoàn toàn mạng Giá trị $V_\psi$.}
      \vspace{0.04cm}
      \begin{enumerate}\setlength{\itemsep}{4pt}
        \item Với mỗi truy vấn $x$, lấy mẫu song song một nhóm $G$ quỹ đạo $\{o_1, o_2, \dots, o_G\} \sim \pi_{\theta_{\text{old}}}(\cdot \mid x)$.
        \item Chấm điểm cả nhóm bằng bộ kiểm chứng RLVR: $\{r_1, r_2, \dots, r_G\}$.
        \item Dùng trung bình nhóm $\mu_r$ làm đường cơ sở $b(x)$ và chuẩn hóa Z-score theo $\sigma_r$.
      \end{enumerate}
    \end{column}
  \end{columns}
\end{frame}

% ----------------------------------------------------------------------------
% SLIDE: SƠ ĐỒ KIẾN TRÚC & LUỒNG DỮ LIỆU CỦA RLVR + GRPO (TOÀN TRANG)
% ----------------------------------------------------------------------------
\begin{frame}{Phần 4.2: Kiến trúc \& Luồng Dữ liệu của Hệ thống RLVR + GRPO}{Tối ưu hóa Chính sách thông qua Lấy mẫu Nhóm $G$ Quỹ đạo mà Không cần Mạng Giá trị $V_\psi$}
  \vspace{0.1cm}
  \begin{center}
    \resizebox{0.96\textwidth}{!}{%
      \begin{tikzpicture}
        \node[data, text width=2.4cm, minimum height=1.2cm] (q) at (0, 0) {
          \textbf{Truy vấn $x$}\\(Toán / Thuật toán)
        };

        \node[trainmodel, text width=3.1cm, minimum height=1.5cm] (policy) at (3.7, 0) {
          \textbf{Chính sách $\pi_\theta$}\\(Đang tối ưu hóa)\\
          \normalfont\scriptsize Lấy mẫu nhóm $G$ quỹ đạo
        };

        \node[frozenmodel, text width=3.1cm, minimum height=0.95cm] (ref) at (3.7, -2.0) {
          \textbf{Tham chiếu $\pi_{\text{ref}}$}\\
          (Điều chuẩn $\DKL$)
        };

        \node[data, text width=2.5cm, minimum height=0.8cm] (o1) at (7.7, 1.35) {Quỹ đạo $o_1$};
        \node[data, text width=2.5cm, minimum height=0.8cm] (o2) at (7.7, 0) {Quỹ đạo $o_2$};
        \node[data, text width=2.5cm, minimum height=0.8cm] (oG) at (7.7, -1.35) {Quỹ đạo $o_G$};

        \node[rewardbox, text width=2.6cm, minimum height=0.8cm] (r1) at (11.2, 1.35) {Kiểm chứng $\to r_1$};
        \node[rewardbox, text width=2.6cm, minimum height=0.8cm] (r2) at (11.2, 0) {Kiểm chứng $\to r_2$};
        \node[rewardbox, text width=2.6cm, minimum height=0.8cm] (rG) at (11.2, -1.35) {Kiểm chứng $\to r_G$};

        \node[lossbox, text width=3.8cm, minimum height=2.6cm] (adv) at (15.6, 0) {
          \textbf{Chuẩn hóa Thống kê Nhóm}\\[4pt]
          $\mu_r = \mean(r_1, \dots, r_G)$\\
          $\sigma_r = \std(r_1, \dots, r_G)$\\[6pt]
          $\hat{A}_i = \dfrac{r_i - \mu_r}{\sigma_r}$
        };

        \draw[flow] (q.east) -- (policy.west);
        \draw[flow] (ref.north) -- (policy.south);

        \coordinate (bus_out) at (5.75, 0);
        \draw[line width=1.05pt, draw=navyblue!85!black] (policy.east) -- (bus_out);
        \draw[flow] (bus_out) |- (o1.west);
        \draw[flow] (bus_out) -- (o2.west);
        \draw[flow] (bus_out) |- (oG.west);

        \draw[flow] (o1.east) -- (r1.west);
        \draw[flow] (o2.east) -- (r2.west);
        \draw[flow] (oG.east) -- (rG.west);

        \draw[flow] (r1.east) -- (13.0, 1.35) |- ([yshift=0.65cm]adv.west);
        \draw[flow] (r2.east) -- (adv.west);
        \draw[flow] (rG.east) -- (13.0, -1.35) |- ([yshift=-0.65cm]adv.west);

        \draw[gradflow] (adv.north) -- (15.6, 2.45) -- node[above, font=\scriptsize\sffamily\bfseries, text=stanfordred] {Cập nhật Đạo hàm $\nabla_\theta \mathcal{J}_{\text{GRPO}}$: Củng cố quỹ đạo $\hat{A}_i > 0$ ($\uparrow$), Chế tài quỹ đạo $\hat{A}_i < 0$ ($\downarrow$)} (3.7, 2.45) -- (policy.north);
      \end{tikzpicture}%
    }
  \end{center}
\end{frame}

% ----------------------------------------------------------------------------
% SLIDE: CÔNG THỨC TOÁN HỌC CỦA GRPO
% ----------------------------------------------------------------------------
\begin{frame}{Phần 4.2: Biểu thức Toán học của Tối ưu hóa Chính sách Tương đối Nhóm (GRPO)}{Hàm Mục tiêu Cắt tỉa Kết hợp Chuẩn hóa Điểm Thưởng Cấp độ Chuỗi (Shao et al., 2024)}
  \begin{block}{Công thức 8: Hàm Mục tiêu GRPO và Lợi thế Chuẩn hóa Tương đối $\hat{A}_i$}
    \footnotesize
    \[
      \mathcal{J}_{\text{GRPO}}(\theta) = \E_{x,\; \{o_i\}_{i=1}^G \sim \pi_{\theta_{\text{old}}}} \left[ \frac{1}{G} \sum_{i=1}^{G} \frac{1}{|o_i|} \sum_{t=1}^{|o_i|} \Big( \min\big(\rho_{i,t}\hat{A}_i,\; \clip(\rho_{i,t}, 1-\epsilon, 1+\epsilon)\hat{A}_i\big) - \beta \DKL(\pi_\theta \parallel \pi_{\text{ref}}) \Big) \right]
    \]
    \vspace{-0.1cm}
    \[
      \text{Với lợi thế chuẩn hóa Z-score cho mọi token } t \in \{1, \dots, |o_i|\}\text{:} \qquad
      \boxed{\hat{A}_i = \frac{r_i - \mean(r_1, r_2, \dots, r_G)}{\std(r_1, r_2, \dots, r_G)}}
    \]
  \end{block}

  \subhead{Đặc tính Giải tích}
  \small
  \begin{itemize}\setlength{\itemsep}{5pt}
    \item \textbf{Đồng nhất lợi thế cấp chuỗi ($\hat{A}_{i,t} \equiv \hat{A}_i$):} Mọi token thuộc quỹ đạo $o_i$ chia sẻ chung trọng số lợi thế $\hat{A}_i$.
    \item \textbf{Tự thích nghi theo độ khó của $x$:} Phép chuẩn hóa Z-score giữ kỳ vọng $\E[\hat{A}_i] = 0$ và phương sai $\mathrm{Var}(\hat{A}_i) = 1$ ổn định trên mọi bài toán.
  \end{itemize}
\end{frame}

% ----------------------------------------------------------------------------
% SLIDE: VÍ DỤ TÍNH TOÁN MINH HỌA GRPO
% ----------------------------------------------------------------------------
\begin{frame}{Phần 4.2: Minh họa Định lượng Cơ chế Phân bổ Lợi thế trong GRPO ($G = 4$)}{Phân tách Tự động giữa Mẫu Dương Giải Đúng và Mẫu Âm Giải Sai qua Chuẩn hóa Nhóm}
  \subhead{Xét truy vấn toán học $x$ với nhóm $G = 4$ quỹ đạo tự sinh $\{o_1, o_2, o_3, o_4\}$ ($\mu_r = 0.275, \sigma_r = 0.42$)}
  \vspace{0.05cm}
  \footnotesize
  \renewcommand{\arraystretch}{1.25}
  \begin{tabularx}{\textwidth}{@{}l l c c X@{}}
    \toprule
    \textbf{Quỹ đạo} & \textbf{Kết quả Kiểm chứng Tất định} & \textbf{Điểm $r_i$} & \textbf{Lợi thế $\hat{A}_i = \frac{r_i - 0.275}{0.42}$} & \textbf{Hướng Cập nhật Đạo hàm} \\
    \midrule
    \textbf{Quỹ đạo $o_1$} & Đáp số \textbf{Đúng} + Đúng thẻ \texttt{<think>} & $r_1 = \mathbf{1.0}$ & $\hat{A}_1 = \mathbf{+1.72 > 0}$ & \textbf{\textcolor{tealgreen}{Tăng mạnh xác suất ($\uparrow\uparrow$)}} \\
    \hline
    \textbf{Quỹ đạo $o_2$} & Đáp số \textbf{Sai}, đúng thẻ \texttt{<think>} & $r_2 = \mathbf{0.1}$ & $\hat{A}_2 = \mathbf{-0.41 < 0}$ & \textbf{\textcolor{stanfordred}{Giảm nhẹ xác suất ($\downarrow$)}} \\
    \hline
    \textbf{Quỹ đạo $o_3$} & Đáp số \textbf{Sai} + Sai cấu trúc & $r_3 = \mathbf{0.0}$ & $\hat{A}_3 = \mathbf{-0.65 < 0}$ & \textbf{\textcolor{stanfordred}{Giảm mạnh xác suất ($\downarrow\downarrow$)}} \\
    \hline
    \textbf{Quỹ đạo $o_4$} & Đáp số \textbf{Sai} + Sai cấu trúc & $r_4 = \mathbf{0.0}$ & $\hat{A}_4 = \mathbf{-0.65 < 0}$ & \textbf{\textcolor{stanfordred}{Giảm mạnh xác suất ($\downarrow\downarrow$)}} \\
    \bottomrule
  \end{tabularx}

  \vspace{0.10cm}
  \begin{exampleblock}{Nhận xét Định lượng}
    \footnotesize
    Không cần mạng Giá trị $V_\psi$, chuẩn hóa nhóm tự động gán lợi thế dương cho quỹ đạo đúng duy nhất ($o_1$) và đẩy các quỹ đạo sai xuống miền lợi thế âm.
  \end{exampleblock}
\end{frame}

% ----------------------------------------------------------------------------
% SLIDE: HAI TINH CHỈNH LÝ THUYẾT CỦA GRPO: DR. GRPO & DAPO (2025)
% ----------------------------------------------------------------------------
\begin{frame}{Phần 4.2: Các Cải tiến Giải tích Thế hệ 2025 đối với GRPO --- Dr.\ GRPO \& DAPO}{Khắc phục Thiên kiến Độ dài Mẫu Âm và Hiện tượng Sụp đổ Entropy Sớm}
  \begin{columns}[T,onlytextwidth]
    \begin{column}{0.48\textwidth}
      \limithead{1. Thiên kiến Độ dài $\longrightarrow$ Dr.\ GRPO (Liu et al., 2025)}
      \footnotesize
      \begin{itemize}\setlength{\itemsep}{5pt}
        \item \textbf{Khiếm khuyết của GRPO gốc:} Chia tổng lợi thế cho độ dài chuỗi $\frac{1}{|o_i|}$ trong Công thức 8.
        \item \textbf{Hệ quả:} Với mẫu âm ($\hat{A}_i < 0$), chuỗi sai càng dài ($|o_i|$ lớn) thì mức phạt trung bình trên mỗi token càng nhẹ $\implies$ mô hình kéo dài câu trả lời sai.
        \item \textbf{Hiệu chỉnh Dr.\ GRPO:} Bỏ hệ số chia $\frac{1}{|o_i|}$ và $\sigma_r$, giữ estimator đạo hàm không chệch.
      \end{itemize}
    \end{column}
    \begin{column}{0.48\textwidth}
      \motivehead{2. Sụp đổ Entropy $\longrightarrow$ DAPO (Yu et al., 2025)}
      \footnotesize
      \begin{itemize}\setlength{\itemsep}{5pt}
        \item \textbf{Hạn chế cắt tỉa đối xứng $[1-\epsilon, 1+\epsilon]$:} Với $\epsilon = 0.2$, token khám phá mới ($\pi_{\text{old}} = 0.01$) bị chặn trên tại $0.012$, kìm hãm đa dạng suy luận.
        \item \textbf{Hiệu chỉnh Clip-Higher (DAPO):} Tách cận cắt tỉa bất đối xứng $\mathbf{[1 - \epsilon_{\text{low}}, \; 1 + \epsilon_{\text{high}}]}$ với $\epsilon_{\text{high}} > \epsilon_{\text{low}}$ để duy trì entropy khám phá.
      \end{itemize}
    \end{column}
  \end{columns}
\end{frame}

\subsection{4.3. Sự Nảy sinh Chuỗi Suy luận Nội tại, Mở rộng lúc Suy diễn \& Agentic RL (2025--2026)}

% ----------------------------------------------------------------------------
% SLIDE: SỰ NẢY SINH TƯ DUY HỆ THỐNG 2 (<think>)
% ----------------------------------------------------------------------------
\begin{frame}{Phần 4.3: Sự Nảy sinh Hành vi Tự Phản biện \& Chuỗi Suy luận Nội tại (\texttt{<think>})}{Phân tích Thực nghiệm trên Mô hình DeepSeek-R1-Zero (DeepSeek-AI, 2025)}
  \begin{columns}[T,onlytextwidth]
    \begin{column}{0.48\textwidth}
      \subhead{Thiết kế Phần thưởng Tối giản}
      \footnotesize
      Chạy GRPO trực tiếp trên Base Model (không qua SFT), chỉ dùng \textbf{2 quy tắc kiểm chứng}:
      \begin{enumerate}\setlength{\itemsep}{5pt}
        \item \textbf{Phần thưởng Độ chính xác ($r_{\text{acc}}$):}\\
          Kiểm tra tính tương đương đại số của đáp số cuối.
        \item \textbf{Phần thưởng Định dạng ($r_{\text{format}}$):}\\
          Kiểm tra cấu trúc thẻ \texttt{<think>...</think>} và \texttt{<answer>...</answer>}.
      \end{enumerate}
      \vspace{0.04cm}
      \scriptsize
      \textit{Chú ý:} Không có bất kỳ hàm thưởng nào quy định nội dung hay độ dài bên trong thẻ \texttt{<think>}.
    \end{column}
    \begin{column}{0.48\textwidth}
      \motivehead{Cơ chế Tự Hình thành Tư duy Hệ thống 2}
      \footnotesize
      \begin{itemize}\setlength{\itemsep}{5pt}
        \item Các quỹ đạo tự phân bổ nhiều token trong \texttt{<think>} để \textbf{phân rã bài toán, kiểm định giả thuyết và tự sửa lỗi} đạt xác suất giải đúng ($r_{\text{acc}} = 1$) vượt trội.
        \item Qua chọn lọc tự nhiên của $\hat{A}_i > 0$, chính sách tự củng cố hành vi siêu nhận thức (\textit{Metacognitive Self-Correction}).
        \item Chuyển dịch từ \textbf{Phản xạ Hệ thống 1} sang \textbf{Suy luận Hệ thống 2}.
      \end{itemize}
    \end{column}
  \end{columns}
\end{frame}

% ----------------------------------------------------------------------------
% SLIDE: MỞ RỘNG TÍNH TOÁN LÚC SUY DIỄN & AGENTIC RL (2025-2026)
% ----------------------------------------------------------------------------
\begin{frame}{Phần 4.3: Mở rộng Tính toán Lúc Suy diễn \& Agentic RL Đa vòng (2025--2026)}{Từ Suy luận Nội tại Đơn bước sang Tác tử Tương tác Môi trường Thực thi}
  \begin{columns}[T,onlytextwidth]
    \begin{column}{0.48\textwidth}
      \motivehead{1. Mở rộng Lúc Suy diễn (Snell et al., 2024)}
      \footnotesize
      Hiệu năng suy luận tăng theo ngân sách tính toán lúc giải mã trên 2 trục trực giao:
      \begin{itemize}\setlength{\itemsep}{5pt}
        \item \textbf{Mở rộng Nối tiếp (Sequential Scaling):}\\
          Kéo dài chuỗi \texttt{<think>} để tự kiểm tra và hiệu chỉnh sai số trên một quỹ đạo.
        \item \textbf{Mở rộng Song song (Parallel Scaling):}\\
          Lấy mẫu $N$ quỹ đạo độc lập, hợp nhất bằng \textbf{Self-Consistency (Majority Voting)} hoặc \textbf{Best-of-$N$}.
      \end{itemize}
    \end{column}
    \begin{column}{0.48\textwidth}
      \motivehead{2. Xu hướng 2025--2026: Agentic \& Rubric RL}
      \footnotesize
      \begin{itemize}\setlength{\itemsep}{5pt}
        \item \textbf{Agentic RL Đa vòng (Multi-turn Tool-Use RL):}\\
          Nhúng GRPO vào môi trường thực thi: suy nghĩ \texttt{<think>} $\to$ gọi công cụ (Terminal, Compiler, Search) $\to$ nhận phản hồi môi trường $\to$ suy luận tiếp.
        \item \textbf{Rubric-based RL \& GenRM:}\\
          Mở rộng RLVR sang tác vụ mở bằng bộ tiêu chí kiểm chứng phân rã (Checklists) và mô hình thưởng có suy luận (\textit{Generative Reward Model}).
      \end{itemize}
    \end{column}
  \end{columns}
\end{frame}

\subsection{4.4. Tổng kết Phần 4 (RLVR + GRPO): Thành tựu \& Các Vấn đề Tồn đọng cho Phần 5}

% ----------------------------------------------------------------------------
% SLIDE: TỔNG KẾT PHẦN 4 (1/2): NHỮNG GÌ ĐÃ ĐẠT ĐƯỢC & TỐI ƯU SO VỚI DPO & PPO
% ----------------------------------------------------------------------------
\begin{frame}{Phần 4.4: Tổng kết Phần 4 (1 / 2) --- Thành tựu Đạt được so với DPO \& PPO}{Đánh giá Đột phá Suy luận Hệ thống 2 của RLVR + GRPO (Cuối 2024--2026)}
  \motivehead{Những gì RLVR + GRPO Đã Giải quyết \& Tối ưu hơn so với Phần 2 và Phần 3}
  \vspace{0.05cm}
  \footnotesize
  \renewcommand{\arraystretch}{1.32}
  \begin{tabularx}{\textwidth}{@{}p{3.2cm} X X@{}}
    \toprule
    \textbf{Tiêu chí Đánh giá} & \textbf{PPO (Phần 2) \& DPO (Phần 3)} & \textbf{RLVR + GRPO (Phần 4 -- 2024--2026)} \\
    \midrule
    \textbf{1. Tín hiệu định giá} & Mạng $r_\phi$ hoặc cặp tĩnh (bị \textit{Reward Hacking} trên Toán/Code) & \textbf{\textcolor{tealgreen}{Bộ kiểm chứng tất định (RLVR)}}: Khách quan tuyệt đối, triệt tiêu gian lận hình thức \\
    \hline
    \textbf{2. Ước lượng đường cơ sở} & PPO cần mạng Critic $V_\psi$ lớn, sai số cao trên chuỗi dài & \textbf{\textcolor{tealgreen}{Chuẩn hóa nhóm $G$ quỹ đạo (GRPO)}}: Khử hoàn toàn mạng Critic $V_\psi$, giảm $50\%$ VRAM \\
    \hline
    \textbf{3. Năng lực suy luận} & Phản xạ Hệ thống 1 (học tĩnh trên dữ liệu sẵn có) & \textbf{\textcolor{tealgreen}{Tư duy Hệ thống 2 (\texttt{<think>})}}: Tự khám phá trực tuyến và mở rộng tính toán lúc chạy \\
    \bottomrule
  \end{tabularx}
\end{frame}

% ----------------------------------------------------------------------------
% SLIDE: TỔNG KẾT PHẦN 4 (2/2): VẤN ĐỀ TỒN ĐỌNG MÀ PHẦN 5 (CHƯNG CẤT) GIẢI QUYẾT
% ----------------------------------------------------------------------------
\begin{frame}{Phần 4.4: Tổng kết Phần 4 (2 / 2) --- Các Vấn đề Tồn đọng Dẫn tới Phần 5 (Chưng cất)}{Tại sao RLVR + GRPO Chỉ Khám phá Hiệu quả trên Mô hình Frontier mà Đình trệ trên Mô hình Nhỏ?}
  \limithead{Ba Rào cản Khi Chạy Trực tiếp RLVR + GRPO trên Mô hình Dung lượng Nhỏ (1.5B--14B)}
  \small
  \begin{enumerate}\setlength{\itemsep}{6pt}
    \item \textbf{Điều kiện Cần của Giải thuật GRPO ($\sigma_r > 0$):}\\
      \footnotesize Trong nhóm $G$ quỹ đạo $\{o_1, \dots, o_G\}$, bắt buộc phải xuất hiện ít nhất một vài mẫu dương giải đúng để tạo độ lệch chuẩn $\sigma_r > 0$.
    \item \textbf{Hiện tượng Cạn kiệt Phần thưởng \& Triệt tiêu Đạo hàm ($\nabla_\theta \mathcal{J} = \mathbf{0}$):}\\
      \footnotesize Trước bài toán khó, mô hình nhỏ (1.5B--14B) giải sai cả $G = 64$ lần thử ($r_1 = \dots = r_G = 0 \implies \hat{A}_i = 0 \implies \mathbf{\nabla_\theta \mathcal{J} = 0}$). Quá trình học hoàn toàn bất động!
    \item \textbf{Chi phí Tính toán của Khám phá RL từ Đầu:}\\
      \footnotesize Chỉ khả thi trên mô hình Frontier dung lượng lớn (như 671B).
  \end{enumerate}

  \vspace{0.06cm}
  \begin{exampleblock}{Bài toán Chuyển giao sang Phần 5 (Chưng cất Tri thức -- Knowledge Distillation)}
    \footnotesize
    Sử dụng mô hình Frontier đã qua RL làm \textbf{Mô hình Thầy ($\pi_T$)} để truyền thẳng phân phối dự đoán $P_T^{(\tau)}$ và quỹ đạo suy luận \texttt{<think>} sang \textbf{Mô hình Trò ($\pi_S$)}.
  \end{exampleblock}
\end{frame}

% ============================================================================
% PHẦN 5: CHƯNG CẤT TRI THỨC, KHÔNG GIAN THAM SỐ & TỔNG KẾT HỆ THỐNG
% ============================================================================
\section{Phần 5: Chưng cất Tri thức, Không gian Tham số \& Tổng kết Hệ thống}

\subsection{5.1. Chưng cất Tri thức từ Mô hình Thầy Frontier sang Mô hình Trò Nhỏ (Distillation)}

% ----------------------------------------------------------------------------
% SLIDE: TỔNG QUAN PHẦN 5.1 — KHUNG THẦY-TRÒ
% ----------------------------------------------------------------------------
\begin{frame}{Phần 5.1: Lý thuyết Chưng cất Tri thức (Knowledge Distillation)}{Khắc phục Hiện tượng Triệt tiêu Đạo hàm RL trên Mô hình Nhỏ (Từ Phần 4.4)}
  \subhead{1. Khuôn khổ Học Thầy--Trò (Hinton et al., 2015; Kim \& Rush, 2016)}
  \small
  Chuyển giao phân phối xác suất và quỹ đạo suy luận từ \textbf{Mô hình Thầy Frontier ($\pi_T$)} đã qua huấn luyện RL quy mô lớn sang \textbf{Mô hình Trò ($\pi_S$)} có dung lượng tham số nhỏ.

  \vspace{0.10cm}
  \motivehead{2. Ba Phương thức Biểu diễn Tri thức Chuyển giao}
  \small
  \begin{enumerate}\setlength{\itemsep}{5pt}
    \item \textbf{Chưng cất Phân phối Dự đoán (Logit / White-box KD):}\\
      \footnotesize Đối khớp phân phối xác suất mềm $P_T^{(\tau)}(\cdot \mid s_t)$ trên toàn bộ từ điển $\mathcal{V}$ tại từng token.
    \item \textbf{Chưng cất Chuỗi Đầu ra (Sequence-Level / Black-box KD):}\\
      \footnotesize Tối đa hóa hợp lý logarit trên phản hồi hoàn chỉnh $y_T \sim \pi_T(\cdot \mid x)$.
    \item \textbf{Chưng cất Quỹ đạo Suy luận (Rationale / Reasoning KD):}\\
      \footnotesize Chuyển giao đồng thời chuỗi lập luận trung gian \texttt{<think>} $z_T$ đã kiểm chứng và kết luận $y_T$.
  \end{enumerate}
\end{frame}

% ----------------------------------------------------------------------------
% SLIDE: BẰNG CHỨNG THỰC NGHIỆM: CHƯNG CẤT VS TỰ CHẠY RL TRÊN MÔ HÌNH NHỎ
% ----------------------------------------------------------------------------
\begin{frame}{Phần 5.1: Đối chứng Thực nghiệm --- Chưng cất Tri thức so với Tự Tối ưu hóa RL}{Phân tích Hiệu năng trên Mô hình Qwen-32B-Base (DeepSeek-AI, 2025)}
  \subhead{Thí nghiệm Đối chứng: Tự Chạy GRPO từ Đầu so với Chưng cất từ Thầy 671B}
  \vspace{0.06cm}
  \footnotesize
  \renewcommand{\arraystretch}{1.28}
  \begin{tabularx}{\textwidth}{@{}p{4.8cm} c c X@{}}
    \toprule
    \textbf{Phương thức Huấn luyện trên Qwen-32B} & \textbf{AIME 2024} & \textbf{MATH-500} & \textbf{Giải thích Cơ chế Lý thuyết} \\
    \midrule
    \textbf{1. Tự chạy RL từ đầu (Zero-RL / GRPO $>10{,}000$ bước)} & $47.0\%$ & $83.9\%$ & Bị giới hạn bởi dung lượng khám phá tiên nghiệm của mạng 32B \\
    \hline
    \textbf{2. Chưng cất Suy luận từ Thầy DeepSeek-R1 671B} & \textbf{\textcolor{tealgreen}{72.6\%}} & \textbf{\textcolor{tealgreen}{94.3\%}} & \textbf{Vượt trội $+25.6\%$} nhờ kế thừa trực tiếp quỹ đạo suy luận của Thầy 671B \\
    \bottomrule
  \end{tabularx}

  \vspace{0.12cm}
  \begin{exampleblock}{Phân công Chiến lược trong Hệ sinh thái Hậu Huấn luyện}
    \footnotesize
    Dùng \textbf{Online RL (RLVR + GRPO)} trên mô hình Frontier để \textbf{khám phá tri thức suy luận mới}; dùng \textbf{Chưng cất Quỹ đạo Suy luận} để nén năng lực đó xuống mô hình nhỏ.
  \end{exampleblock}
\end{frame}

% ----------------------------------------------------------------------------
% SLIDE: SƠ ĐỒ BA HÌNH THỨC CHUYỂN GIAO TRI THỨC TRONG CHƯNG CẤT (TOÀN TRANG)
% ----------------------------------------------------------------------------
\begin{frame}{Phần 5.1: Kiến trúc \& Luồng Dữ liệu của Ba Phương thức Chưng cất Tri thức}{So sánh Luồng Tín hiệu Huấn luyện giữa Chưng cất Hộp Trắng (White-box) và Hộp Đen (Black-box)}
  \vspace{0.15cm}
  \begin{center}
    \resizebox{0.96\textwidth}{!}{%
      \begin{tikzpicture}
        \node[frozenmodel, text width=3.3cm, minimum height=2.9cm, fill=softblue, draw=navyblue, line width=1.1pt] (teacher) at (0, 0) {
          \textbf{Mô hình Thầy $\pi_T$}\\
          (Frontier Teacher)\\[6pt]
          \normalfont\scriptsize Đã qua Pre-training,\\
          SFT \& RL (GRPO)\\
          (Ví dụ: 671B tham số)
        };

        \node[rewardbox, text width=5.6cm, minimum height=1.05cm] (k1) at (6.4, 1.5) {
          \textbf{1. Chưng cất Phân phối Dự đoán (Logit / White-box KD)}\\
          \normalfont\scriptsize Chuyển giao phân phối mềm $P_T^{(\tau)}(v \mid x, y_{<t})$ trên toàn bộ từ điển $\mathcal{V}$
        };

        \node[data, text width=5.6cm, minimum height=1.05cm] (k2) at (6.4, 0) {
          \textbf{2. Chưng cất Chuỗi Đầu ra (Sequence / Black-box KD)}\\
          \normalfont\scriptsize Chuyển giao phản hồi hoàn chỉnh $y_T \sim \pi_T(\cdot \mid x)$ làm chuỗi mục tiêu
        };

        \node[lossbox, text width=5.6cm, minimum height=1.05cm, fill=softgreen, draw=tealgreen] (k3) at (6.4, -1.5) {
          \textbf{3. Chưng cất Quỹ đạo Suy luận (Rationale / Reasoning KD)}\\
          \normalfont\scriptsize Chuyển giao đồng thời chuỗi suy luận $z_T$ (\texttt{<think>}) và đáp án $y_T$ đã kiểm chứng
        };

        \node[trainmodel, text width=3.3cm, minimum height=2.9cm] (student) at (12.8, 0) {
          \textbf{Mô hình Trò $\pi_S$}\\
          (Student Model)\\[6pt]
          \normalfont\scriptsize Dung lượng nhỏ gọn,\\
          tối ưu hóa độ trễ\\
          (Ví dụ: 1.5B -- 32B)
        };

        \coordinate (bus_left) at (2.4, 0);
        \draw[line width=1.05pt, draw=navyblue!85!black] (teacher.east) -- (bus_left);
        \draw[flow] (bus_left) |- (k1.west);
        \draw[flow] (bus_left) -- (k2.west);
        \draw[flow] (bus_left) |- (k3.west);

        \draw[flow] (k1.east) -- (10.4, 1.5) |- ([yshift=0.7cm]student.west);
        \draw[flow] (k2.east) -- (student.west);
        \draw[flow] (k3.east) -- (10.4, -1.5) |- ([yshift=-0.7cm]student.west);
      \end{tikzpicture}%
    }
  \end{center}
  \vspace{0.05cm}
  \footnotesize
  \textbf{Điều kiện kỹ thuật:} Phương thức 1 đòi hỏi cùng bộ từ vựng $\mathcal{V}$ và truy cập vector logit (\textit{White-box}); Phương thức 2 \& 3 chỉ dùng văn bản đầu ra nên áp dụng được cho mọi cặp mô hình (\textit{Black-box}).
\end{frame}

% ----------------------------------------------------------------------------
% SLIDE: PHƯƠNG THỨC 1 — CHƯNG CẤT DỰ ĐOÁN (LOGIT KD) & NHIỆT ĐỘ TAU
% ----------------------------------------------------------------------------
\begin{frame}{Phần 5.1: Phương thức 1 --- Chưng cất Phân phối Dự đoán (Logit Distillation)}{Biểu thức Softmax Nhiệt độ $\tau$ và Cấu trúc Tương đồng Ngữ nghĩa Nội tại (Hinton et al., 2015)}
  \begin{block}{Công thức 9: Phân phối Softmax Nhiệt độ $\tau$ và Hàm Mục tiêu Logit KD}
    \[
      P^{(\tau)}(v_k \mid s_t) = \frac{\exp(z_k / \tau)}{\sum_{j=1}^{|\mathcal{V}|} \exp(z_j / \tau)}, \qquad
      \mathcal{L}_{\text{Logit-KD}}(\theta_S) = \tau^2 \, \DKL\Big( P_T^{(\tau)}(\cdot \mid s_t) \;\Big\|\; P_S^{(\tau)}(\cdot \mid s_t) \Big)
    \]
  \end{block}

  \begin{columns}[T,onlytextwidth]
    \begin{column}{0.50\textwidth}
      \motivehead{Cấu trúc Tương quan Ngữ nghĩa (Dark Knowledge)}
      \footnotesize
      \begin{itemize}\setlength{\itemsep}{4pt}
        \item Nhãn cứng (One-hot) chỉ gán xác suất $1$ cho 1 từ và $0$ cho $|\mathcal{V}|-1$ từ còn lại.
        \item Phân phối mềm $P_T^{(\tau)}$ (ví dụ: \textit{``mèo''}: $0.70$, \textit{``chó''}: $0.25$, \textit{``tủ lạnh''}: $10^{-8}$) truyền tải \textbf{cấu trúc hình học tương đồng giữa các khái niệm}.
      \end{itemize}
    \end{column}
    \begin{column}{0.46\textwidth}
      \subhead{Vai trò của Nhiệt độ $\tau > 1$ và Hệ số $\tau^2$}
      \footnotesize
      \begin{itemize}\setlength{\itemsep}{4pt}
        \item \textbf{Chia cho $\tau > 1$:} Làm phẳng phân phối Softmax, nâng xác suất vùng đuôi để Trò tiếp nhận rõ cấu trúc tương quan.
        \item \textbf{Nhân với $\tau^2$:} Bù lại sự suy giảm biên độ đạo hàm $\nabla_{z_S}\DKL \propto 1/\tau^2$.
      \end{itemize}
    \end{column}
  \end{columns}
\end{frame}

% ----------------------------------------------------------------------------
% SLIDE: PHƯƠNG THỨC 2 & 3 — SEQUENCE KD VS RATIONALE KD
% ----------------------------------------------------------------------------
\begin{frame}{Phần 5.1: Phương thức 2 \& 3 --- Sequence KD so với Rationale Distillation}{Khắc phục Sự Đứt gãy Biểu diễn Trung gian bằng Chuỗi Suy luận Nội tại \texttt{<think>}}
  \begin{block}{Hàm Mục tiêu Chưng cất Quỹ đạo Suy luận (Rationale Distillation -- Hsieh et al., 2023)}
    \[
      \mathcal{L}_{\text{Rationale-KD}}(\theta_S) = -\E_{(x, z_T, y_T) \sim \mathcal{D}_{\text{verified}}} \Big[ \log \pi_S(z_T \mid x) + \log \pi_S(y_T \mid x, z_T) \Big]
    \]
  \end{block}

  \begin{columns}[T,onlytextwidth]
    \begin{column}{0.48\textwidth}
      \limithead{Hạn chế của Sequence KD ($x \mapsto y_T$)}
      \footnotesize
      \begin{itemize}\setlength{\itemsep}{5pt}
        \item Đáp án cuối $y_T$ là kết quả của nhiều bước biến đổi ẩn bên trong mạng Thầy lớn.
        \item Lược bỏ chuỗi suy luận trung gian $z_T$ khiến Mô hình Trò (ít tầng ẩn hơn) rơi vào trạng thái \textbf{khớp mẫu bề mặt (Surface Pattern Matching)}.
      \end{itemize}
    \end{column}
    \begin{column}{0.48\textwidth}
      \motivehead{Ưu thế của Rationale KD ($x \mapsto z_T \mapsto y_T$)}
      \footnotesize
      \begin{itemize}\setlength{\itemsep}{5pt}
        \item Lọc mẫu theo ngưỡng kiểm chứng (\textit{Rejection Sampling}) để giữ các cặp $(z_T, y_T)$ đúng của Thầy.
        \item Sinh chuỗi \texttt{<think>} $z_T$ trước giúp phân rã bài toán khó thành chuỗi các bước dự đoán cục bộ đơn giản trên ngữ cảnh.
      \end{itemize}
    \end{column}
  \end{columns}
\end{frame}

% ----------------------------------------------------------------------------
% SLIDE: ĐỐI CHIẾU FORWARD KL VS REVERSE KL TRONG CHƯNG CẤT
% ----------------------------------------------------------------------------
\begin{frame}{Phần 5.1: Toán học Chưng cất Hiện đại --- Forward KL so với Reverse KL}{Tại sao Reverse KL (Mode-Seeking) Tối ưu cho Mô hình Trò Dung lượng Nhỏ? (Gu et al., 2024)}
  \begin{block}{Công thức 10: Hai Chiều Phân kỳ KL giữa Mô hình Thầy ($P_T$) và Mô hình Trò ($P_S$)}
    \[
      \text{\textbf{1. Forward KL:}} \; \DKL(P_T \parallel P_S) = \sum_{y} P_T(y) \log \frac{P_T(y)}{P_S(y)}
      \qquad
      \text{\textbf{2. Reverse KL:}} \; \DKL(P_S \parallel P_T) = \sum_{y} P_S(y) \log \frac{P_S(y)}{P_T(y)}
    \]
  \end{block}

  \begin{columns}[T,onlytextwidth]
    \begin{column}{0.48\textwidth}
      \limithead{Forward KL: Phủ Trung bình (Mean-Seeking)}
      \footnotesize
      \begin{itemize}\setlength{\itemsep}{4pt}
        \item Khi $P_T(y) > 0$, bắt buộc $P_S(y) > 0$ (nếu $P_S \to 0 \implies \log\frac{P_T}{P_S} \to +\infty$).
        \item Mô hình Trò nhỏ cố trải rộng phân phối bao phủ mọi đỉnh của Thầy $\implies$ trọng tâm rơi vào vùng mật độ thấp ở giữa (nơi $P_T(y) \approx 0$).
      \end{itemize}
    \end{column}
    \begin{column}{0.48\textwidth}
      \motivehead{Reverse KL: Hội tụ Cực đại (Mode-Seeking)}
      \footnotesize
      \begin{itemize}\setlength{\itemsep}{4pt}
        \item Khi $P_T(y) = 0$, bắt buộc $P_S(y) = 0$. Nhưng nếu Trò bỏ qua một số đỉnh phụ của Thầy ($P_S(y) = 0$), tích số $0 \log 0 = 0$ không gây phạt.
        \item Cho phép Trò \textbf{tập trung $100\%$ dung lượng vào một cực đại tốt nhất của Thầy}.
      \end{itemize}
    \end{column}
  \end{columns}
\end{frame}

% ----------------------------------------------------------------------------
% SLIDE: QUY TRÌNH KẾT HỢP TRÊN MÔ HÌNH NHỎ: DISTILLATION -> RL / DPO
% ----------------------------------------------------------------------------
\begin{frame}{Phần 5.1: Quy trình Kết hợp trên Mô hình Nhỏ (Distillation $\longrightarrow$ GRPO / DPO)}{Tại sao Mô hình Nhỏ Hiện đại Phối hợp Cả Chưng cất Lẫn Học Tăng cường?}
  \subhead{Quy trình Hai Pha Khắc phục Đồng thời Triệt tiêu Đạo hàm và Exposure Bias}
  \small
  \begin{enumerate}\setlength{\itemsep}{7pt}
    \item \textbf{Pha 1 -- Chưng cất từ Mô hình Thầy Frontier (Logit KD / Rationale KD):}\\
      \footnotesize Khởi tạo năng lực suy luận cho mô hình nhỏ, nâng xác suất giải đúng tiên nghiệm từ $0\%$ lên $50\%\text{--}70\%$. Tuy nhiên, do học tĩnh trên chuỗi của Thầy, mô hình Trò vẫn chịu \textit{Exposure Bias} và thiếu tín hiệu từ mẫu âm của chính nó.
    \item \textbf{Pha 2 -- Chạy Bồi GRPO (RLVR) và DPO Trực tiếp trên Mô hình Trò:}\\
      \footnotesize Nhờ Pha 1, nhóm $G$ quỹ đạo lấy mẫu trên mô hình nhỏ đã có cả mẫu dương lẫn mẫu âm ($\sigma_r > 0 \implies \hat{A}_i \neq 0$). Chạy tiếp GRPO + DPO giúp mô hình nhỏ tự khử lỗi lập luận nội tại và vượt trần chưng cất tĩnh (ứng dụng trên \textit{Qwen3, Gemma 2/3, DeepScaleR}).
  \end{enumerate}
\end{frame}

\subsection{5.2. Thích nghi Hạng Thấp (LoRA), Hợp nhất Tham số \& Đánh giá Hệ thống}

% ----------------------------------------------------------------------------
% SLIDE: LÝ THUYẾT KHÔNG GIAN THAM SỐ: LORA & MODEL MERGING
% ----------------------------------------------------------------------------
\begin{frame}{Phần 5.2: Không gian Tham số --- Thích nghi Hạng Thấp (LoRA) \& Hợp nhất Mô hình}{Tối ưu hóa trên Đa tạp Chiều Thấp và Tổ hợp Tuyến tính Trọng số}
  \begin{columns}[T,onlytextwidth]
    \begin{column}{0.48\textwidth}
      \motivehead{1. Chiều Nội tại Thấp \& LoRA (Hu et al., 2022)}
      \footnotesize
      \[
        W = W_0 + \Delta W = W_0 + \frac{\alpha}{r} B A, \quad r \ll \min(d, k)
      \]
      \begin{itemize}\setlength{\itemsep}{4pt}
        \item \textbf{Chiều nội tại thấp (Aghajanyan et al.):} Độ lệch $\Delta W$ trong Post-Training nằm trên đa tạp con hạng $r$ rất nhỏ ($r \in \{16, 32\}$).
        \item \textbf{Chống Quên Thảm khốc (Biderman et al., 2024):} Đóng băng $W_0$ và chỉ cập nhật $<1\%$ tham số giúp bảo toàn tri thức gốc của Base Model.
      \end{itemize}
    \end{column}
    \begin{column}{0.48\textwidth}
      \motivehead{2. Số học Vector Tác vụ \& Hợp nhất Mô hình}
      \footnotesize
      \[
        \tau_i = \theta_{\text{finetuned}}^{(i)} - \theta_{\text{base}} \implies \theta_{\text{merged}} = \theta_{\text{base}} + \sum_{i=1}^{M} \lambda_i \tau_i
      \]
      \begin{itemize}\setlength{\itemsep}{4pt}
        \item \textbf{Linear Mode Connectivity:} Các mô hình tinh chỉnh từ cùng $\theta_{\text{base}}$ nằm chung trong một vùng trũng hàm mất mát (Loss Basin).
        \item \textbf{TIES-Merging \& DARE:} Cắt tỉa $90\%$ biến thiên nhỏ và đồng bộ dấu để hợp nhất kỹ năng \textit{Toán} ($\tau_{\text{math}}$) và \textit{Hội thoại} ($\tau_{\text{chat}}$).
      \end{itemize}
    \end{column}
  \end{columns}
\end{frame}

% ----------------------------------------------------------------------------
% SLIDE: LÝ THUYẾT ĐÁNH GIÁ HỆ THỐNG: ALIGNMENT TAX, ELO & LLM-AS-A-JUDGE
% ----------------------------------------------------------------------------
\begin{frame}{Phần 5.2: Lý thuyết Đánh giá Hệ thống --- Thuế Căn chỉnh, Elo \& Trọng tài AI}{Đo lường Sự Đánh đổi Hiệu năng và Hiệu chỉnh Thiên kiến Đánh giá Tự động (Zheng et al., 2023)}
  \begin{columns}[T,onlytextwidth]
    \begin{column}{0.45\textwidth}
      \limithead{1. Thuế Căn chỉnh \& Hệ thống Elo}
      \scriptsize
      \begin{itemize}\setlength{\itemsep}{4pt}
        \item \textbf{Thuế Căn chỉnh (Alignment Tax):}\\
          Quá trình RLHF / DPO thu hẹp entropy phân phối, làm giảm nhẹ độ đa dạng hoặc độ hiệu chỉnh xác suất (Calibration).
        \item \textbf{Định lượng Elo (Chatbot Arena):}\\
          Đánh giá đối kháng mù từng cặp theo mô hình Bradley-Terry:
          \[
            P(A \succ B) = \frac{1}{1 + 10^{(R_B - R_A)/400}}
          \]
      \end{itemize}
    \end{column}
    \begin{column}{0.52\textwidth}
      \motivehead{2. Ba Thiên kiến của Trọng tài AI (LLM-as-a-Judge)}
      \scriptsize
      \renewcommand{\arraystretch}{1.22}
      \begin{tabularx}{\textwidth}{@{}l X@{}}
        \toprule
        \textbf{Thiên kiến} & \textbf{Phương pháp Hiệu chỉnh Thống kê} \\
        \midrule
        \textbf{1. Vị trí} (\textit{Position}) & Đánh giá 2 chiều hoán vị $(A, B)$ và $(B, A)$ \\
        \hline
        \textbf{2. Độ dài} (\textit{Verbosity}) & Hồi quy tách phương sai độ dài (\textit{Length-Controlled Win Rate}) \\
        \hline
        \textbf{3. Tự tương đồng} & Hội đồng đa mô hình hoặc bộ kiểm chứng tất định \\
        \bottomrule
      \end{tabularx}
    \end{column}
  \end{columns}
\end{frame}

\subsection{5.3. Tổng kết Toàn diện Chuyên đề Seminar (2020--2026)}

% ----------------------------------------------------------------------------
% SLIDE: BẢN ĐỒ PHÂN LOẠI 3 TẦNG THUẬT NGỮ (TOÀN TRANG)
% ----------------------------------------------------------------------------
\begin{frame}{Phần 5.3: Hệ thống Phân loại Học thuật Ba Tầng trong Hậu Huấn luyện}{Tách biệt Rạch ròi giữa Nguồn Phản hồi, Độ mịn Gán Tín dụng và Giải thuật Tối ưu hóa}
  \vspace{0.1cm}
  \begin{center}
    \resizebox{0.95\textwidth}{!}{%
      \begin{tikzpicture}[
        layer/.style={rectangle, rounded corners=4pt, draw=navyblue, text=white, text width=3.3cm, minimum height=1.25cm, align=center, font=\scriptsize\sffamily\bfseries},
        item/.style={rectangle, rounded corners=4pt, draw=tealgreen!85!black, line width=0.9pt, fill=softgreen, text width=3.5cm, minimum height=1.25cm, align=center, font=\scriptsize\sffamily}
      ]
        \node[layer, fill=navyblue] (L1) at (0, 1.65) {TẦNG 1: Nguồn Phát\\Tín hiệu Phản hồi\\(Feedback Source)};
        \node[item] (L1a) at (4.1, 1.65) {\textbf{RLHF (Phần 2)}\\Đánh giá từ Con người\\(Cặp đối chiếu $y_w \succ y_l$)};
        \node[item] (L1b) at (8.1, 1.65) {\textbf{RLAIF (Phần 3)}\\Đánh giá từ Mô hình AI\\(Hiến pháp / Quy tắc)};
        \node[item, fill=softyellow, draw=amberorange!90!black] (L1c) at (12.1, 1.65) {\textbf{RLVR (Phần 4)}\\Bộ Kiểm chứng Tất định\\(Toán học / Trình biên dịch)};

        \node[layer, fill=tealgreen] (L2) at (0, 0) {TẦNG 2: Độ Mịn\\Gán Tín dụng\\(Credit Granularity)};
        \node[item, text width=5.5cm] (L2a) at (5.1, 0) {\textbf{ORM (Outcome Reward Model)}\\Cấp phát phần thưởng duy nhất ở cuối chuỗi phản hồi};
        \node[item, text width=5.5cm] (L2b) at (11.1, 0) {\textbf{PRM (Process Reward Model)}\\Cấp phát phần thưởng chi tiết sau từng bước suy luận trung gian};

        \node[layer, fill=stanfordred] (L3) at (0, -1.65) {TẦNG 3: Giải thuật\\Tối ưu hóa Tham số\\(Optimization Algorithm)};
        \node[item, fill=softred, draw=stanfordred] (L3a) at (4.1, -1.65) {\textbf{PPO (Phần 2)}\\RL trực tuyến có mạng Giá trị $V_\psi$\\(Hệ thống 4 mô hình)};
        \node[item, fill=softred, draw=stanfordred] (L3b) at (8.1, -1.65) {\textbf{DPO / SimPO (Phần 3)}\\Tối ưu hóa trực tiếp từ cặp tĩnh\\(Nghiệm đóng hàm mục tiêu)};
        \node[item, fill=softred, draw=stanfordred] (L3c) at (12.1, -1.65) {\textbf{GRPO (Phần 4)}\\RL trực tuyến chuẩn hóa nhóm $G$\\(Loại bỏ mạng Giá trị)};
      \end{tikzpicture}%
    }
  \end{center}
  \vspace{0.05cm}
  \footnotesize
  \textbf{Tổ hợp kiến trúc điển hình theo lịch sử:} \textit{InstructGPT (2022)} = \textbf{RLHF} + \textbf{ORM} + \textbf{PPO}. \quad $|$ \quad \textit{DeepSeek-R1 (2025)} = \textbf{RLVR} + \textbf{ORM} + \textbf{GRPO}.
\end{frame}

% ----------------------------------------------------------------------------
% SLIDE: SƠ ĐỒ TOÀN CẢNH CHUỖI TIẾN HÓA LÝ THUYẾT POST-TRAINING (TOÀN TRANG)
% ----------------------------------------------------------------------------
\begin{frame}{Phần 5.3: Sơ đồ Tổng hợp Chuỗi Tiến hóa Lý thuyết (2020--2026)}{Mạch Logic Nhân quả: Giới hạn của Phương pháp Tiền nhiệm $\longrightarrow$ Động cơ của Phương pháp Kế tiếp}
  \vspace{0.15cm}
  \begin{center}
    \resizebox{0.96\textwidth}{!}{%
      \begin{tikzpicture}[
        method/.style={rectangle, rounded corners=4pt, draw=navyblue, line width=1.05pt, fill=softblue, text width=2.9cm, minimum height=1.2cm, align=center, font=\scriptsize\sffamily\bfseries},
        limit/.style={rectangle, rounded corners=4pt, draw=stanfordred, line width=0.9pt, fill=softred, text width=3.0cm, minimum height=1.1cm, align=center, font=\tiny\sffamily}
      ]
        \node[method] (pre) at (0, 1.65) {2020: Pre-training\\$\mathcal{L}_{\text{MLE}}(u_t \mid u_{<t})$};
        \node[limit] (lim1) at (3.6, 1.65) {\textbf{Giới hạn:}\\Lệch mục tiêu với tác vụ\\tuân thủ chỉ thị};
        \node[method] (sft) at (7.2, 1.65) {Phần 1: SFT (2021)\\$\mathcal{L}_{\text{SFT}}(y^* \mid x)$};
        \node[limit] (lim2) at (10.8, 1.65) {\textbf{Giới hạn:}\\Lệch phơi nhiễm, ảo giác,\\thiếu tín hiệu từ mẫu âm};

        \node[method] (rlhf) at (10.8, 0) {Phần 2: RLHF + PPO\\(2022 -- 4 mô hình)};
        \node[limit] (lim3) at (7.2, 0) {\textbf{Giới hạn:}\\Chi phí bộ nhớ 4 mạng,\\sinh mẫu chậm, bất ổn định};
        \node[method, fill=softgreen, draw=tealgreen] (dpo) at (3.6, 0) {Phần 3: DPO (2023)\\(Nghiệm đóng khử $Z(x)$)};
        \node[limit] (lim4) at (0, 0) {\textbf{Giới hạn DPO:}\\Học tĩnh, thất bại ở\\Toán/Code (\textit{Reward Hacking})};

        \node[method, fill=softyellow, draw=amberorange!90!black] (grpo) at (0, -1.65) {Phần 4: RLVR + GRPO\\(2024--2026: Bỏ Critic $V_\psi$)};
        \node[limit] (lim5) at (5.4, -1.65) {\textbf{Giới hạn trên Mô hình nhỏ:}\\Cạn kiệt phần thưởng ($\hat{A}_i = 0$),\\đạo hàm triệt tiêu ($\nabla_\theta = \mathbf{0}$)};
        \node[method, fill=softgreen, draw=tealgreen] (kd) at (10.8, -1.65) {Phần 5: Chưng cất\\(Logit / Sequence /\\Rationale KD)};

        \draw[flow] (pre.east) -- (lim1.west);
        \draw[flow] (lim1.east) -- (sft.west);
        \draw[flow] (sft.east) -- (lim2.west);
        \draw[flow] (lim2.south) -- (rlhf.north);
        \draw[flow] (rlhf.west) -- (lim3.east);
        \draw[flow] (lim3.west) -- (dpo.east);
        \draw[flow] (dpo.west) -- (lim4.east);
        \draw[flow] (lim4.south) -- (grpo.north);
        \draw[flow] (grpo.east) -- (lim5.west);
        \draw[flow] (lim5.east) -- (kd.west);
      \end{tikzpicture}%
    }
  \end{center}
  \vspace{0.05cm}
  \footnotesize
  \textbf{Quy trình tích hợp hiện đại:} \textbf{SFT khởi tạo} $\longrightarrow$ \textbf{RLVR + GRPO khám phá suy luận trên Frontier} $\longrightarrow$ \textbf{DPO căn chỉnh đa mục tiêu} $\longrightarrow$ \textbf{Distillation nén xuống mô hình nhỏ}.
\end{frame}

% ----------------------------------------------------------------------------
% SLIDE: BẢNG ĐỐI CHIẾU TỔNG KẾT 5 PHẦN TRỤ CỘT POST-TRAINING
% ----------------------------------------------------------------------------
\begin{frame}{Phần 5.3: Bảng Tổng hợp Đối chiếu Năm Phần Trụ cột theo Trình tự Lịch sử}{Tóm tắt Thành tựu Đạt được, Cấu hình Hệ thống và Điểm Nghẽn Chuyển giao qua Từng Giai đoạn}
  \vspace{0.02cm}
  \scriptsize
  \renewcommand{\arraystretch}{1.18}
  \begin{tabularx}{\textwidth}{@{}p{2.35cm} p{2.15cm} p{2.05cm} X X@{}}
    \toprule
    \textbf{Giai đoạn / Phần} & \textbf{Tín hiệu Học} & \textbf{Số Mô hình} & \textbf{Thành tựu Tối ưu so với Trước đó} & \textbf{Vấn đề Để lại cho Phần Sau} \\
    \midrule
    \textbf{Phần 1: SFT}\\(2020--2021) & Mẫu dương chuẩn $(x, y^*)$ & \textbf{1 mạng:} $\pi_\theta$ & Giải quyết \textit{Objective Mismatch}, dạy khuôn mẫu hội thoại nhanh & \textit{Exposure Bias}, kích thích ảo giác, thiếu tín hiệu từ mẫu âm \\
    \hline
    \textbf{Phần 2: RLHF+PPO}\\(2022 -- ChatGPT) & Điểm thưởng $r_\phi(x, \hat{y})$ trên mẫu tự sinh & \textbf{4 mạng:} $\pi_\theta, \pi_{\text{ref}}, r_\phi, V_\psi$ & Học trực tuyến (\textit{On-policy}), đạo hàm 2 chiều chế tài mẫu âm & Tốn bộ nhớ gấp 4 lần, sinh mẫu chậm, đạo hàm RL bất ổn định \\
    \hline
    \textbf{Phần 3: DPO}\\(2023--2024) & Cặp mẫu dương/âm $(x, y_w \succ y_l)$ & \textbf{2 mạng:} $\pi_\theta, \pi_{\text{ref}}$ & Nghiệm đóng khử $Z(x)$, bỏ $r_\phi$ \& $V_\psi$, ổn định như SFT & Học tĩnh (\textit{Off-policy}), thất bại ở Toán/Code do \textit{Reward Hacking} \\
    \hline
    \textbf{Phần 4: RLVR+GRPO}\\(Cuối 2024--2026) & Bộ kiểm chứng trên nhóm $\{o_1..o_G\}$ & \textbf{2 mạng:} $\pi_\theta, \pi_{\text{ref}}$ (+ Verifier) & Bỏ Critic $V_\psi$, miễn nhiễm \textit{Reward Hacking}, nảy sinh \texttt{<think>} & Mô hình nhỏ bị cạn kiệt phần thưởng ($\hat{A}_i = 0 \implies \nabla_\theta = \mathbf{0}$) \\
    \hline
    \textbf{Phần 5: Chưng cất}\\(Distillation) & Phân phối $P_T^{(\tau)}$ hoặc chuỗi $(z_T, y_T)$ & \textbf{2 mạng:} Thầy $\pi_T \to$ Trò $\pi_S$ & Truyền thẳng suy luận từ Thầy Frontier xuống mô hình nhỏ & Kết hợp thêm GRPO/DPO nhẹ sau chưng cất để khử \textit{Exposure Bias} \\
    \bottomrule
  \end{tabularx}
\end{frame}

% ----------------------------------------------------------------------------
% SLIDE: KẾT LUẬN CHUYÊN ĐỀ SEMINAR & THẢO LUẬN HỌC THUẬT
% ----------------------------------------------------------------------------
\begin{frame}{Phần 5.3: Tổng kết Chuyên đề Seminar \& Phiên Thảo luận Học thuật}{Năm Nguyên lý Lý thuyết Trọng tâm Xuyên suốt Lịch sử Hậu Huấn luyện (2020--2026)}
  \subhead{Các Kết luận Khoa học Cốt lõi}
  \scriptsize
  \begin{enumerate}\setlength{\itemsep}{5pt}
    \item \textbf{Phần 1 (2020--2021 -- SFT):} Định hình phân phối tuân thủ chỉ thị $\pi_\theta(y^* \mid x)$, nhưng chịu giới hạn bởi \textit{Exposure Bias} và sự vắng mặt của tín hiệu phạt từ mẫu âm.
    \item \textbf{Phần 2 (2022 -- RLHF \& PPO):} Kết hợp hàm thưởng Bradley-Terry $r_\phi(y_w \succ y_l)$, điều chuẩn vùng tin cậy $\beta \DKL$ và giải thuật PPO giúp căn chỉnh trực tuyến, nhưng chịu gánh nặng bộ nhớ 4 mô hình.
    \item \textbf{Phần 3 (2023--2024 -- DPO):} Nghiệm đóng phân phối Gibbs--Boltzmann kết hợp tính bất biến tịnh tiến Bradley-Terry giúp khử hàm phân hoạch $\beta \log Z(x)$, tối ưu hóa trực tiếp trên cặp mẫu dương/âm $(y_w, y_l)$.
    \item \textbf{Phần 4 (Cuối 2024--2026 -- RLVR \& GRPO):} Bộ kiểm chứng tất định kết hợp chuẩn hóa nhóm $G$ quỹ đạo giúp xóa bỏ mạng Giá trị $V_\psi$ và kích hoạt chuỗi suy luận nội tại Hệ thống 2 (\texttt{<think>}).
    \item \textbf{Phần 5 (Chưng cất Tri thức \& Không gian Tham số):} Chưng cất quỹ đạo suy luận kết hợp phân kỳ KL ngược (\textit{Reverse KL}) giúp mô hình nhỏ vượt qua hiện tượng triệt tiêu đạo hàm của RL.
  \end{enumerate}

  \vspace{0.08cm}
  \begin{center}
    \normalsize \textbf{\textcolor{navyblue}{Trân trọng cảm ơn sự theo dõi của Quý Thầy Cô và các bạn! --- Thảo luận \& Hỏi Đáp (Q\&A)}}
  \end{center}
\end{frame}


\end{document}
